
\documentclass[10pt,a4paper]{ctexart}
\usepackage{geometry}
\geometry{left=1.4cm,right=1.4cm,top=1.15cm,bottom=1.0cm,headsep=0.3cm,footskip=0.5cm}
\usepackage{fontspec}
\setmainfont{Times New Roman}
\setsansfont{Helvetica}
\usepackage{xcolor}
\usepackage{booktabs}
\usepackage{tabularx}
\usepackage{array}
\usepackage{amssymb}
\xeCJKDeclareCharClass{CJK}{"2460 -> "247F}
\usepackage{fancyhdr}
\usepackage[hidelinks]{hyperref}
\definecolor{accent}{HTML}{0E4D5C}
\definecolor{warn}{HTML}{9B2226}
\definecolor{gray}{HTML}{5A5A5A}
\definecolor{rule}{HTML}{CBDCDD}
\definecolor{soft}{HTML}{F2F7F7}
\definecolor{ink}{HTML}{1C1C1E}
\renewcommand{\familydefault}{\rmdefault}
\setlength{\parindent}{0pt}
\setlength{\parskip}{0pt}
\linespread{1.0}
\pagestyle{fancy}
\fancyhf{}
\renewcommand{\headrulewidth}{0pt}
\fancyhead[L]{\footnotesize\color{gray}英语二小作文 · 十二类文体模板}
\fancyhead[R]{\footnotesize\color{gray}Part A · 应用文}
\fancyfoot[C]{\footnotesize\color{gray}\thepage}
\renewcommand{\arraystretch}{1.08}

\begin{document}
\clearpage
\thispagestyle{empty}
\vspace*{1.0cm}
{\color{accent}\footnotesize\bfseries 考研英语（二）Part A · 应用文写作}\par
\vspace{8pt}
{\color{accent}\rule{\linewidth}{2.2pt}}\par\vspace{8pt}
{\Huge\bfseries 英语二小作文 · 十二类文体模板}\par
\vspace{6pt}
{\large\color{gray} 十二类文体 · 每类一页 · 骨架可替换 · 按 10 分档写法组织}\par
\vspace{8pt}
{\color{accent}\rule{\linewidth}{0.8pt}}\par\vspace{12pt}
\par\vspace{1.5pt}{\small 覆盖 2010—2025 全部真题文体：邀请 / 回复 / 建议 / 感谢 / 祝贺 / 道歉 / 投诉 / 咨询 / 通知 / 汇报}
\par\vspace{1.5pt}{\small 补充大纲内两类：推荐信、申请信（模拟题高频）}
\par\vspace{1.5pt}{\small 每页固定五块：题型识别 → 结构骨架 → 可替换句型 → 满分范文 → 雷区与得分点}
\par\vspace{1.5pt}{\small 范文均按“正文约 100 词”控制，可直接背诵，也可按【槽位】改写}
\vspace{14pt}
\noindent{\color{accent}\bfseries\large 目\quad 录}\par\vspace{3pt}
{\color{rule}\rule{\linewidth}{0.6pt}}\par\vspace{3pt}
\noindent\begin{tabularx}{\linewidth}{@{}Xr@{}}
\footnotesize 01\quad 邀请信 & \footnotesize\hfill \pageref{g01} \\
\footnotesize 02\quad 回复邀请（接受 · 婉拒） & \footnotesize\hfill \pageref{g02} \\
\footnotesize 03\quad 建议信 & \footnotesize\hfill \pageref{g03} \\
\footnotesize 04\quad 感谢信 & \footnotesize\hfill \pageref{g04} \\
\footnotesize 05\quad 祝贺信 & \footnotesize\hfill \pageref{g05} \\
\footnotesize 06\quad 道歉信 & \footnotesize\hfill \pageref{g06} \\
\footnotesize 07\quad 投诉信 & \footnotesize\hfill \pageref{g07} \\
\footnotesize 08\quad 咨询信 & \footnotesize\hfill \pageref{g08} \\
\footnotesize 09\quad 通知 & \footnotesize\hfill \pageref{g09} \\
\footnotesize 10\quad 汇报 · 告知信 & \footnotesize\hfill \pageref{g10} \\
\footnotesize 11\quad 推荐信 & \footnotesize\hfill \pageref{g11} \\
\footnotesize 12\quad 申请 · 求职信 & \footnotesize\hfill \pageref{g12} \\
\end{tabularx}\par
\vfill
\noindent\colorbox{soft}{\parbox{\dimexpr\linewidth-2\fboxsep\relax}{%
\footnotesize\color{gray}编写依据：考研英语（二）考试大纲对 Part A 的要求，以及工作区《批改指令》中的评分档位与批改原则。全部英文为原创句型，与市面通用“万能模板”不重合。}}\par
\vspace{4pt}
\noindent{\footnotesize\color{gray}使用建议：先看第 02 页的决策表定文体，再翻到对应页背“骨架句”，最后用第 03 页的六条自查过一遍。}
\clearpage
\label{rules}
\vspace*{-6pt}
{\color{accent}\bfseries\LARGE 通用规则：先拿满格式分，再谈语言分}\par\vspace{4pt}
{\color{accent}\rule{\linewidth}{1.2pt}}\par
\vspace{4pt}\noindent{\color{accent}\bfseries 一、分数构成（10 分）}\par\nopagebreak
{\color{rule}\rule{\linewidth}{0.6pt}}\par\nopagebreak
\vspace{1pt}
\vspace{2pt}\noindent\begingroup\footnotesize\renewcommand{\arraystretch}{1.15}
\begin{tabularx}{\dimexpr\linewidth-1.4pt\relax}{@{}>{\raggedright\arraybackslash}p{3.2cm}>{\raggedright\arraybackslash}p{3.0cm}>{\raggedright\arraybackslash}X@{}}
\toprule
{\color{accent}\bfseries 评分维度} & {\color{accent}\bfseries 权重感受} & {\color{accent}\bfseries 最容易失分的地方} \\
\midrule
任务完成 & 最高 & 漏写时间/地点/联系方式；合并题只答一半 \\
语言准确 & 高 & 主谓一致、时态混用、搭配错（congratulations for） \\
格式语气 & 中 & 称呼与署名不合身份；公务信写成口语便条 \\
连贯与变化 & 中 & 通篇 I will / you should 单一句式；无衔接手段 \\
\bottomrule
\end{tabularx}\endgroup\par
\vspace{2pt}
\vspace{4pt}\noindent{\color{accent}\bfseries 二、格式四要素（书信）}\par\nopagebreak
{\color{rule}\rule{\linewidth}{0.6pt}}\par\nopagebreak
\vspace{1pt}
\vspace{2pt}\noindent\begingroup\footnotesize\renewcommand{\arraystretch}{1.15}
\begin{tabularx}{\dimexpr\linewidth-1.4pt\relax}{@{}>{\raggedright\arraybackslash}p{3.2cm}>{\raggedright\arraybackslash}p{6.4cm}>{\raggedright\arraybackslash}X@{}}
\toprule
{\color{accent}\bfseries 位置} & {\color{accent}\bfseries 写法} & {\color{accent}\bfseries 例} \\
\midrule
称呼 & 顶格，逗号结尾；公务用姓，私人用名 & Dear Professor Wang, / Dear Jack, \\
正文 & 2—3 段，每段一个功能，段首不空格 & 见各文体骨架 \\
结束语 & 顶格，逗号结尾，与署名同一语气层 & Yours sincerely, / Best wishes, \\
署名 & 另起一行，题目指定姓名 & Li Ming \\
\bottomrule
\end{tabularx}\endgroup\par
\vspace{2pt}
\vspace{4pt}\noindent{\color{accent}\bfseries 三、通知类格式（无称呼）}\par\nopagebreak
{\color{rule}\rule{\linewidth}{0.6pt}}\par\nopagebreak
\vspace{1pt}
\vspace{2pt}\noindent\begingroup\footnotesize\renewcommand{\arraystretch}{1.15}
\begin{tabularx}{\dimexpr\linewidth-1.4pt\relax}{@{}>{\raggedright\arraybackslash}p{3.2cm}>{\raggedright\arraybackslash}p{6.4cm}>{\raggedright\arraybackslash}X@{}}
\toprule
{\color{accent}\bfseries 位置} & {\color{accent}\bfseries 写法} & {\color{accent}\bfseries 例} \\
\midrule
标题 & 居中，单独一行，首字母大写 & Notice \\
正文 & 1—3 段：目的 → 安排 → 号召/联系方式 & 见第 09 页 \\
落款 & 右下角两行：单位 + 日期 & The Student Union / Oct. 20, 2025 \\
\bottomrule
\end{tabularx}\endgroup\par
\vspace{2pt}
\vspace{4pt}\noindent{\color{accent}\bfseries 四、词数与时间}\par\nopagebreak
{\color{rule}\rule{\linewidth}{0.6pt}}\par\nopagebreak
\vspace{1pt}
\vspace{1pt}\begin{list}{}{\leftmargin=1.3em\labelwidth=0.95em\labelsep=0.35em\itemindent=0pt\topsep=0pt\parsep=1.5pt\itemsep=1pt}
\item[{\color{accent}\scriptsize$\blacktriangleright$}]{\footnotesize 正文约 100 词（不含称呼、结束语、署名）。120 词以内不必担心，超过 150 词会挤占大作文时间并放大错误。}
\item[{\color{accent}\scriptsize$\blacktriangleright$}]{\footnotesize 考场建议 18 分钟：审题画要点 3 分钟 → 写作 13 分钟 → 检查 2 分钟（主谓一致、时态、拼写、格式）。}
\item[{\color{accent}\scriptsize$\blacktriangleright$}]{\footnotesize 三段最稳：第一段点题、第二段展开、第三段收束。四段也可以，但每段不能只有一句话。}
\end{list}\vspace{2pt}
\vspace{4pt}\noindent{\color{accent}\bfseries 五、评分档位（据《批改指令》，不人为封顶）}\par\nopagebreak
{\color{rule}\rule{\linewidth}{0.6pt}}\par\nopagebreak
\vspace{1pt}
\vspace{2pt}\noindent\begingroup\footnotesize\renewcommand{\arraystretch}{1.15}
\begin{tabularx}{\dimexpr\linewidth-1.4pt\relax}{@{}>{\raggedright\arraybackslash}p{2.2cm}>{\raggedright\arraybackslash}X@{}}
\toprule
{\color{accent}\bfseries 档位} & {\color{accent}\bfseries 表现} \\
\midrule
8—9 分 & 很好完成任务：要点完整，语言准确自然，表达有变化，行文连贯，格式与语气恰当 \\
6—7 分 & 较好完成任务：主要要点完整，允许遗漏少量次要信息；语言基本准确，偶有错误 \\
4—5 分 & 基本完成任务：涵盖多数要点；有一些语言错误，但不影响整体理解 \\
2—3 分 & 遗漏重要要点或夹杂无关内容；语言错误较多，影响理解 \\
1 分 / 0 分 & 未完成任务，语言错误严重 / 未作答或内容完全无关 \\
\bottomrule
\end{tabularx}\endgroup\par
\vspace{2pt}
\vspace{4pt}\noindent{\color{accent}\bfseries 六、六个高频扣分点（自查）}\par\nopagebreak
{\color{rule}\rule{\linewidth}{0.6pt}}\par\nopagebreak
\vspace{1pt}
\vspace{1pt}\begingroup\setlength{\parindent}{0pt}
\begin{list}{}{\leftmargin=1.3em\labelwidth=0.95em\labelsep=0.35em\itemindent=0pt\topsep=0pt\partopsep=0pt\parsep=1.5pt\itemsep=0pt}
\item[{\color{warn}\scriptsize$\blacksquare$}]{\footnotesize 漏要点：邀请信不写时间地点、投诉信不写购买渠道、咨询信不列具体问题。}
\item[{\color{warn}\scriptsize$\blacksquare$}]{\footnotesize 称呼错：Dear Mr. Jack（名不加 Mr.）、Dear Teacher（不这么说，用 Dear Professor Wang）。}
\item[{\color{warn}\scriptsize$\blacksquare$}]{\footnotesize 语气错：投诉信全程 angry、trash，或建议信通篇 you must。}
\item[{\color{warn}\scriptsize$\blacksquare$}]{\footnotesize 搭配错：congratulations for → on；grateful to your help → for；depend of → on。}
\item[{\color{warn}\scriptsize$\blacksquare$}]{\footnotesize 时态乱：邀请用将来时，感谢用过去时，汇报用现在完成时；同一段内不要来回跳。}
\item[{\color{warn}\scriptsize$\blacksquare$}]{\footnotesize 字数失衡：正文只有 60 词，或把 Mid-Autumn Festival 展开成三句话 cultural background。}
\end{list}
\endgroup
\clearpage
\label{decide}
{\color{accent}\bfseries\LARGE 看到题干，先归到哪一类}\par\vspace{4pt}
{\color{accent}\rule{\linewidth}{1.2pt}}\par
\vspace{4pt}\noindent{\footnotesize 英语二 Part A 的题干通常给两件事：写什么（事务）+ 要你完成什么动作（邀请、说明、请求）。先把“动作词”找出来，文体就定了；再看有没有“合并任务”，合并题必须两个动作都落地。}\par\vspace{4pt}
\vspace{2pt}\noindent\begingroup\footnotesize\renewcommand{\arraystretch}{1.15}
\begin{tabularx}{\dimexpr\linewidth-1.4pt\relax}{@{}>{\raggedright\arraybackslash}p{7.6cm}>{\raggedright\arraybackslash}p{3.6cm}>{\raggedright\arraybackslash}p{1.6cm}>{\raggedright\arraybackslash}X@{}}
\toprule
{\color{accent}\bfseries 题干里的信号词} & {\color{accent}\bfseries 文体} & {\color{accent}\bfseries 页码} & {\color{accent}\bfseries 最容易漏的要点} \\
\midrule
invite / 邀请……参加…… & 邀请信 & 01 & 时间地点 + 请对方确认 \\
reply to / 回复……的来信 & 回复信（接受·婉拒） & 02 & 第一段必须有明确态度 \\
advice / suggestions / 就……提建议 & 建议信 & 03 & 建议要带理由，两条不重复 \\
thank / grateful / 感谢 & 感谢信 & 04 & 帮助的具体内容与影响 \\
congratulate / 祝贺 & 祝贺信 & 05 & 祝贺什么 + 为何值得 \\
apologize / 为……道歉 & 道歉信 & 06 & 原因一句 + 补救方案 \\
complain / 投诉 & 投诉信 & 07 & 购买渠道 + 明确要求与期限 \\
information / ask about / 询问 & 咨询信 & 08 & 2—3 个具体问题 \\
Notice / 通知……参加…… & 通知 & 09 & 标题 + 落款 + 联系方式 \\
report / tell sb. about / 汇报 & 汇报·告知信 & 10 & 时间节点与分工 \\
recommend / 推荐 & 推荐信 & 11 & 关系交代 + 具体事例 \\
apply for / 申请 & 申请·求职信 & 12 & 信息来源 + 可量化成果 \\
\bottomrule
\end{tabularx}\endgroup\par
\vspace{2pt}
\vspace{4pt}\noindent{\color{accent}\bfseries 三条通用判断}\par\nopagebreak
{\color{rule}\rule{\linewidth}{0.6pt}}\par\nopagebreak
\vspace{1pt}
\vspace{1pt}\begin{list}{}{\leftmargin=1.35em\labelwidth=1.0em\labelsep=0.35em\itemindent=0pt\topsep=0pt\parsep=2pt\itemsep=2pt}
\item[{\color{accent}\footnotesize\bfseries 1.}]{\footnotesize 合并题（如 2010 感谢＋邀请、2011 祝贺＋建议、2014 告知＋咨询、2019/2024 汇报＋建议）：两件事各写一段，段首用 As for… / Besides… 明确切换，阅卷人一眼能看到两个要点。}
\item[{\color{accent}\footnotesize\bfseries 2.}]{\footnotesize 题干给了姓名、职位、日期、地点，一律照抄进正文，不要自创；题干没给的时间地点，用 next Saturday / Room 203 这类具体而安全的信息补齐。}
\item[{\color{accent}\footnotesize\bfseries 3.}]{\footnotesize 落款永远用题目指定姓名（通常是 Li Ming），不要写自己的名字。}
\end{list}
\clearpage
\clearpage
\label{g01}
\noindent\colorbox{accent}{%
\parbox{\dimexpr\linewidth-2\fboxsep\relax}{%
\color{white}\vspace{1pt}
{\LARGE\bfseries 01}\hspace{0.5em}{\Large 邀请信}\hfill{\large\itshape Invitation}
\\[1pt]{\footnotesize 真题 2010 · 2021 · 2022 · 2025}\par\vspace{1pt}}}
\par
\vspace{4pt}\noindent\colorbox{soft}{%
\parbox{\dimexpr\linewidth-2\fboxsep\relax}{\footnotesize{\color{accent}\bfseries 识别信号}\quad 题干出现 invite / 邀请……参加……，并要求说明活动安排；信尾往往再要一句“请对方确认能否出席”。}}
\par
\vspace{2pt}
\vspace{4pt}\noindent{\color{accent}\bfseries 一、任务要点}\par\nopagebreak
{\color{rule}\rule{\linewidth}{0.6pt}}\par\nopagebreak
\vspace{1pt}
\vspace{2pt}\noindent\begin{tabularx}{\dimexpr\linewidth-1.4pt\relax}{@{}>{\raggedright\arraybackslash}p{2.1cm}X@{}}
\toprule
\footnotesize{\color{accent}\bfseries 写作身份} & \footnotesize 活动组织者、同班同学或社团负责人；用题目给的姓名落款 \\
\midrule
\footnotesize{\color{accent}\bfseries 必须写到} & \footnotesize （1）活动名称与目的　（2）时间、地点　（3）具体安排或亮点　（4）请对方在某日期前回复 \\
\midrule
\footnotesize{\color{accent}\bfseries 一句话策略} & \footnotesize 第一段只点题不铺陈，第二段全部用来交代“什么时间、在哪、做什么”，第三段促回复 \\
\bottomrule
\end{tabularx}\par
\vspace{2pt}
\vspace{4pt}\noindent{\color{accent}\bfseries 二、结构骨架}\par\nopagebreak
{\color{rule}\rule{\linewidth}{0.6pt}}\par\nopagebreak
\vspace{1pt}
\vspace{2pt}\noindent\begingroup\footnotesize\renewcommand{\arraystretch}{1.15}
\begin{tabularx}{\dimexpr\linewidth-1.4pt\relax}{@{}>{\raggedright\arraybackslash}p{2.4cm}>{\raggedright\arraybackslash}p{3.4cm}>{\raggedright\arraybackslash}X@{}}
\toprule
{\color{accent}\bfseries 位置} & {\color{accent}\bfseries 功能} & {\color{accent}\bfseries 必须落到纸面的内容} \\
\midrule
称呼 & 定关系与语气 & 不熟用 Dear Professor Wang,；熟人用 Dear Jack, \\
第一段 & 点题 & 写信目的 + 活动名称（一句话，不写背景） \\
第二段 & 交代安排 & 时间、地点、两个环节或亮点、对方能得到什么 \\
第三段 & 促成回复 & 希望出席 + 请在某日期前确认 / 说明报名方式 \\
落款 & 格式 & Yours sincerely, 换行顶格写 Li Ming \\
\bottomrule
\end{tabularx}\endgroup\par
\vspace{2pt}
\vspace{4pt}\noindent{\color{accent}\bfseries 三、可替换骨架句}\par\nopagebreak
{\color{rule}\rule{\linewidth}{0.6pt}}\par\nopagebreak
\vspace{1pt}
\vspace{1pt}\begingroup\setlength{\parindent}{0pt}
\vspace{2pt}{\footnotesize\bfseries\color{ink}（1） 开头点题}\par\vspace{1pt}
\begin{list}{}{\leftmargin=1.3em\labelwidth=0.95em\labelsep=0.35em\itemindent=0pt\topsep=0pt\partopsep=0pt\parsep=1.5pt\itemsep=0pt}
\item[{\color{accent}\scriptsize$\blacktriangleright$}]{\footnotesize I am writing to invite you to our {\color{accent}\bfseries 【活动名称】}, which will be held on {\color{accent}\bfseries 【日期】}.}\par\vspace{0.5pt}{\footnotesize\color{gray}最稳的开头：which 从句顺手交代时间，一句完成两件事。}
\item[{\color{accent}\scriptsize$\blacktriangleright$}]{\footnotesize On behalf of {\color{accent}\bfseries 【主办方】}, I would like to invite you to {\color{accent}\bfseries 【活动名称】}.}\par\vspace{0.5pt}{\footnotesize\color{gray}代表集体发出邀请，语气更正式，适合公务信。}
\end{list}
\vspace{2pt}{\footnotesize\bfseries\color{ink}（2） 交代安排}\par\vspace{1pt}
\begin{list}{}{\leftmargin=1.3em\labelwidth=0.95em\labelsep=0.35em\itemindent=0pt\topsep=0pt\partopsep=0pt\parsep=1.5pt\itemsep=0pt}
\item[{\color{accent}\scriptsize$\blacktriangleright$}]{\footnotesize The event will start at {\color{accent}\bfseries 【时间】} in {\color{accent}\bfseries 【地点】} and last for about {\color{accent}\bfseries 【时长】}.}\par\vspace{0.5pt}{\footnotesize\color{gray}时间地点一句交代完，避免连着两句都以 I will 开头。}
\item[{\color{accent}\scriptsize$\blacktriangleright$}]{\footnotesize It will begin with {\color{accent}\bfseries 【环节一】}, followed by {\color{accent}\bfseries 【环节二】}.}\par\vspace{0.5pt}{\footnotesize\color{gray}程序用 begin with / followed by，比 first… second… 紧凑。}
\end{list}
\vspace{2pt}{\footnotesize\bfseries\color{ink}（3） 促成回复}\par\vspace{1pt}
\begin{list}{}{\leftmargin=1.3em\labelwidth=0.95em\labelsep=0.35em\itemindent=0pt\topsep=0pt\partopsep=0pt\parsep=1.5pt\itemsep=0pt}
\item[{\color{accent}\scriptsize$\blacktriangleright$}]{\footnotesize I do hope you can make it. Please let me know before {\color{accent}\bfseries 【日期】} whether you are available.}\par\vspace{0.5pt}{\footnotesize\color{gray}make it 意为“能来”，口语但得体，do 用来加强语气。}
\item[{\color{accent}\scriptsize$\blacktriangleright$}]{\footnotesize I would be grateful if you could confirm your attendance by {\color{accent}\bfseries 【日期】}.}\par\vspace{0.5pt}{\footnotesize\color{gray}公务邀请优先用这句，confirm attendance 是标准搭配。}
\end{list}
\vspace{1pt}{\footnotesize\color{gray}语气与落款：公务与师长用前者；朋友用后者。两者后面都要另起一行顶格署名。}\par
\endgroup
\vspace{4pt}\noindent{\color{accent}\bfseries 四、满分范文}\par\nopagebreak
{\color{rule}\rule{\linewidth}{0.6pt}}\par\nopagebreak
\vspace{1pt}
\vspace{2pt}\noindent\colorbox{soft}{\parbox{\dimexpr\linewidth-2\fboxsep\relax}{%
\vspace{2pt}
{\footnotesize\bfseries\color{accent}2022 真题话题：邀请国际学生参加校园美食节}\hfill{\footnotesize\color{gray}正文 105 词}
\par\vspace{2pt}
\footnotesize Dear Jack,\par
\footnotesize I am writing to invite you to our Campus Food Festival, which will take place next Saturday on the main square.\par
\footnotesize The festival will open at 10 a.m. and last until 4 p.m. Students from more than ten countries will set up stalls and offer typical dishes from their hometowns. There will also be a short cooking contest in the afternoon, and I have signed us up as a team.\par
\footnotesize I know you are fond of Chinese food, so I am sure you will enjoy it. Please let me know by Thursday whether you can come, so that I can reserve a place for you.\par
\footnotesize Yours sincerely,\par
\footnotesize Li Ming\par
\par\vspace{2pt}{\footnotesize\color{gray}亲爱的杰克：我写信邀请你参加我们的校园美食节，活动将于下周六在主广场举行。美食节上午十点开始，持续到下午四点。来自十多个国家的学生将摆摊，提供各自家乡的特色菜。下午还有一场简短的厨艺比赛，我已经替我们俩报名组队了。我知道你喜欢中国菜，相信你会喜欢这个活动。请在周四前告诉我能否前来，我好为你留位置。你真诚的，李明}
\par\vspace{1pt}}}\par
\vspace{2pt}
\vspace{4pt}\noindent{\color{accent}\bfseries 五、雷区与得分点}\par\nopagebreak
{\color{rule}\rule{\linewidth}{0.6pt}}\par\nopagebreak
\vspace{1pt}
\vspace{1pt}\begingroup\setlength{\parindent}{0pt}
\begin{list}{}{\leftmargin=1.3em\labelwidth=0.95em\labelsep=0.35em\itemindent=0pt\topsep=0pt\partopsep=0pt\parsep=1.5pt\itemsep=0pt}
\item[{\color{warn}\scriptsize$\blacksquare$}]{\footnotesize 只写“欢迎参加”却不写时间地点，属于要点缺失，直接掉到 5 分档；时间地点必须落在第二段。}
\item[{\color{warn}\scriptsize$\blacksquare$}]{\footnotesize 称呼不加头衔于名：Dear Mr. Jack 错，Dear Jack 或 Dear Mr. Smith 才对。}
\item[{\color{warn}\scriptsize$\blacksquare$}]{\footnotesize 结尾只写 I hope you can come，没有请对方确认，任务里的“征求回复”就没完成。}
\item[{\color{warn}\scriptsize$\blacksquare$}]{\footnotesize 得分点：活动名首字母大写、环节具体（stalls / cooking contest）、给对方一个来的理由。}
\end{list}
\endgroup
\clearpage
\label{g02}
\noindent\colorbox{accent}{%
\parbox{\dimexpr\linewidth-2\fboxsep\relax}{%
\color{white}\vspace{1pt}
{\LARGE\bfseries 02}\hspace{0.5em}{\Large 回复邀请（接受 · 婉拒）}\hfill{\large\itshape Reply to an Invitation}
\\[1pt]{\footnotesize 真题 2017}\par\vspace{1pt}}}
\par
\vspace{4pt}\noindent\colorbox{soft}{%
\parbox{\dimexpr\linewidth-2\fboxsep\relax}{\footnotesize{\color{accent}\bfseries 识别信号}\quad 题干出现 reply to / 回复某封来信，通常含两个动作：先表态（接受或婉拒），再交代安排或理由。}}
\par
\vspace{2pt}
\vspace{4pt}\noindent{\color{accent}\bfseries 一、任务要点}\par\nopagebreak
{\color{rule}\rule{\linewidth}{0.6pt}}\par\nopagebreak
\vspace{1pt}
\vspace{2pt}\noindent\begin{tabularx}{\dimexpr\linewidth-1.4pt\relax}{@{}>{\raggedright\arraybackslash}p{2.1cm}X@{}}
\toprule
\footnotesize{\color{accent}\bfseries 写作身份} & \footnotesize 受邀者，回信给邀请人；沿用对方的称呼与身份 \\
\midrule
\footnotesize{\color{accent}\bfseries 必须写到} & \footnotesize （1）感谢邀请并明确态度　（2）接受则写到达时间与准备内容，婉拒则写冲突原因与替代方案　（3）收尾致谢 \\
\midrule
\footnotesize{\color{accent}\bfseries 一句话策略} & \footnotesize 第一段必须出现态度词（accept / decline），这是本题的“得分开关” \\
\bottomrule
\end{tabularx}\par
\vspace{2pt}
\vspace{4pt}\noindent{\color{accent}\bfseries 二、结构骨架}\par\nopagebreak
{\color{rule}\rule{\linewidth}{0.6pt}}\par\nopagebreak
\vspace{1pt}
\vspace{2pt}\noindent\begingroup\footnotesize\renewcommand{\arraystretch}{1.15}
\begin{tabularx}{\dimexpr\linewidth-1.4pt\relax}{@{}>{\raggedright\arraybackslash}p{2.4cm}>{\raggedright\arraybackslash}p{3.4cm}>{\raggedright\arraybackslash}X@{}}
\toprule
{\color{accent}\bfseries 位置} & {\color{accent}\bfseries 功能} & {\color{accent}\bfseries 必须落到纸面的内容} \\
\midrule
称呼 & 沿用邀请人身份 & Dear Professor Smith, \\
第一段 & 感谢 + 表态 & Thank you for your invitation… I am delighted to accept it. \\
第二段 & 落地 & 接受：到达时间、讲什么、需要对方配合什么；婉拒：原因 + 另约时间 \\
第三段 & 收尾 & 期待见面 / 希望改期，再谢一次 \\
落款 & 格式 & Yours sincerely, + Li Ming \\
\bottomrule
\end{tabularx}\endgroup\par
\vspace{2pt}
\vspace{4pt}\noindent{\color{accent}\bfseries 三、可替换骨架句}\par\nopagebreak
{\color{rule}\rule{\linewidth}{0.6pt}}\par\nopagebreak
\vspace{1pt}
\vspace{1pt}\begingroup\setlength{\parindent}{0pt}
\vspace{2pt}{\footnotesize\bfseries\color{ink}（1） 感谢 + 表态}\par\vspace{1pt}
\begin{list}{}{\leftmargin=1.3em\labelwidth=0.95em\labelsep=0.35em\itemindent=0pt\topsep=0pt\partopsep=0pt\parsep=1.5pt\itemsep=0pt}
\item[{\color{accent}\scriptsize$\blacktriangleright$}]{\footnotesize Thank you for your kind invitation to {\color{accent}\bfseries 【活动】}. I am delighted to accept it.}\par\vspace{0.5pt}{\footnotesize\color{gray}一句感谢、一句表态，态度词不能省，也别拖到第二段。}
\item[{\color{accent}\scriptsize$\blacktriangleright$}]{\footnotesize I am afraid I shall have to decline your invitation, much to my regret.}\par\vspace{0.5pt}{\footnotesize\color{gray}婉拒先致歉，much to my regret 表示“很遗憾”，比 I can't go 得体。}
\end{list}
\vspace{2pt}{\footnotesize\bfseries\color{ink}（2） 说明安排 / 理由}\par\vspace{1pt}
\begin{list}{}{\leftmargin=1.3em\labelwidth=0.95em\labelsep=0.35em\itemindent=0pt\topsep=0pt\partopsep=0pt\parsep=1.5pt\itemsep=0pt}
\item[{\color{accent}\scriptsize$\blacktriangleright$}]{\footnotesize I will arrive at {\color{accent}\bfseries 【地点】} twenty minutes early to {\color{accent}\bfseries 【准备事项】}.}\par\vspace{0.5pt}{\footnotesize\color{gray}主动交代细节，阅卷人据此判断你是否真的“接受了”。}
\item[{\color{accent}\scriptsize$\blacktriangleright$}]{\footnotesize As you suggested, I will focus on {\color{accent}\bfseries 【内容】} and leave time for questions.}\par\vspace{0.5pt}{\footnotesize\color{gray}呼应对方要求，是 8 分以上的加分句。}
\end{list}
\vspace{2pt}{\footnotesize\bfseries\color{ink}（3） 收尾}\par\vspace{1pt}
\begin{list}{}{\leftmargin=1.3em\labelwidth=0.95em\labelsep=0.35em\itemindent=0pt\topsep=0pt\partopsep=0pt\parsep=1.5pt\itemsep=0pt}
\item[{\color{accent}\scriptsize$\blacktriangleright$}]{\footnotesize I am looking forward to meeting you and your students.}\par\vspace{0.5pt}{\footnotesize\color{gray}接受场合用，具体到人比 look forward to your reply 自然。}
\item[{\color{accent}\scriptsize$\blacktriangleright$}]{\footnotesize I hope we can find another time that suits you.}\par\vspace{0.5pt}{\footnotesize\color{gray}婉拒场合必须给替代方案，不能只拒绝。}
\end{list}
\vspace{1pt}{\footnotesize\color{gray}语气与落款：给教授、上级用前者；给很正式的长辈可用后者。都要另起一行署名。}\par
\endgroup
\vspace{4pt}\noindent{\color{accent}\bfseries 四、满分范文}\par\nopagebreak
{\color{rule}\rule{\linewidth}{0.6pt}}\par\nopagebreak
\vspace{1pt}
\vspace{2pt}\noindent\colorbox{soft}{\parbox{\dimexpr\linewidth-2\fboxsep\relax}{%
\vspace{2pt}
{\footnotesize\bfseries\color{accent}2017 真题话题：接受邀请并介绍讲座要点}\hfill{\footnotesize\color{gray}正文 88 词}
\par\vspace{2pt}
\footnotesize Dear Professor Smith,\par
\footnotesize Thank you for your kind invitation to give a lecture on translation to our students. I am delighted to accept it.\par
\footnotesize As you suggested, I will focus on the practical skills of translating news texts and leave half an hour for questions. I plan to arrive at the lecture hall twenty minutes early to check the equipment. If you need my slides in advance, please let me know. I will also prepare a one-page handout for each student.\par
\footnotesize I am looking forward to meeting you and your students.\par
\footnotesize Yours sincerely,\par
\footnotesize Li Ming\par
\par\vspace{2pt}{\footnotesize\color{gray}尊敬的史密斯教授：感谢您邀请我为学生做翻译讲座，我很乐意接受。按照您的建议，我将重点讲新闻文本翻译的实用技巧，并留出半小时提问。我打算提前二十分钟到报告厅检查设备。如果您需要我提前提供幻灯片，请告诉我。我还会为每位学生准备一页讲义。期待与您和您的学生见面。您真诚的，李明}
\par\vspace{1pt}}}\par
\vspace{2pt}
\vspace{4pt}\noindent{\color{accent}\bfseries 五、雷区与得分点}\par\nopagebreak
{\color{rule}\rule{\linewidth}{0.6pt}}\par\nopagebreak
\vspace{1pt}
\vspace{1pt}\begingroup\setlength{\parindent}{0pt}
\begin{list}{}{\leftmargin=1.3em\labelwidth=0.95em\labelsep=0.35em\itemindent=0pt\topsep=0pt\partopsep=0pt\parsep=1.5pt\itemsep=0pt}
\item[{\color{warn}\scriptsize$\blacksquare$}]{\footnotesize 表态藏在句中（I think I can come maybe），语气含糊，直接判要点不完整。}
\item[{\color{warn}\scriptsize$\blacksquare$}]{\footnotesize 接受却不写任何具体安排，整封信只有“好的我去”，内容单薄，6 分档。}
\item[{\color{warn}\scriptsize$\blacksquare$}]{\footnotesize 婉拒只写 I am busy，没有替代时间，属于“只拒绝不解决”。}
\item[{\color{warn}\scriptsize$\blacksquare$}]{\footnotesize 得分点：第一段态度明确 + 第二段有可核实的细节（提前到、留提问时间）+ 收尾得体。}
\end{list}
\endgroup
\clearpage
\label{g03}
\noindent\colorbox{accent}{%
\parbox{\dimexpr\linewidth-2\fboxsep\relax}{%
\color{white}\vspace{1pt}
{\LARGE\bfseries 03}\hspace{0.5em}{\Large 建议信}\hfill{\large\itshape Suggestion / Advice}
\\[1pt]{\footnotesize 真题 2011 · 2016 · 2023}\par\vspace{1pt}}}
\par
\vspace{4pt}\noindent\colorbox{soft}{%
\parbox{\dimexpr\linewidth-2\fboxsep\relax}{\footnotesize{\color{accent}\bfseries 识别信号}\quad 题干出现 advice / suggestions / 就……提出建议；常与“祝贺”“回复”合并考，合并题必须两个动作都落地。}}
\par
\vspace{2pt}
\vspace{4pt}\noindent{\color{accent}\bfseries 一、任务要点}\par\nopagebreak
{\color{rule}\rule{\linewidth}{0.6pt}}\par\nopagebreak
\vspace{1pt}
\vspace{2pt}\noindent\begin{tabularx}{\dimexpr\linewidth-1.4pt\relax}{@{}>{\raggedright\arraybackslash}p{2.1cm}X@{}}
\toprule
\footnotesize{\color{accent}\bfseries 写作身份} & \footnotesize 朋友、晚辈或学生，建议要留余地；不要写成上级布置任务 \\
\midrule
\footnotesize{\color{accent}\bfseries 必须写到} & \footnotesize （1）承接对方的好消息或请求　（2）两条建议 + 各自理由　（3）希望采纳并祝好 \\
\midrule
\footnotesize{\color{accent}\bfseries 一句话策略} & \footnotesize 两条建议必须“不同类”：一条讲做什么，一条讲怎么做或注意什么，避免“多读书/多练习”重复 \\
\bottomrule
\end{tabularx}\par
\vspace{2pt}
\vspace{4pt}\noindent{\color{accent}\bfseries 二、结构骨架}\par\nopagebreak
{\color{rule}\rule{\linewidth}{0.6pt}}\par\nopagebreak
\vspace{1pt}
\vspace{2pt}\noindent\begingroup\footnotesize\renewcommand{\arraystretch}{1.15}
\begin{tabularx}{\dimexpr\linewidth-1.4pt\relax}{@{}>{\raggedright\arraybackslash}p{2.4cm}>{\raggedright\arraybackslash}p{3.4cm}>{\raggedright\arraybackslash}X@{}}
\toprule
{\color{accent}\bfseries 位置} & {\color{accent}\bfseries 功能} & {\color{accent}\bfseries 必须落到纸面的内容} \\
\midrule
称呼 & 按关系定 & Dear Mike, / Dear Professor Wang, \\
第一段 & 承接 + 点题 & 听说……／你问我……，我给几点建议 \\
第二段 & 建议 + 理由 & 建议一（最要紧的）＋理由；建议二＋理由或具体做法 \\
第三段 & 收尾 & 希望有用 + 祝顺利，或表示愿意继续帮忙 \\
落款 & 格式 & Yours sincerely, + Li Ming \\
\bottomrule
\end{tabularx}\endgroup\par
\vspace{2pt}
\vspace{4pt}\noindent{\color{accent}\bfseries 三、可替换骨架句}\par\nopagebreak
{\color{rule}\rule{\linewidth}{0.6pt}}\par\nopagebreak
\vspace{1pt}
\vspace{1pt}\begingroup\setlength{\parindent}{0pt}
\vspace{2pt}{\footnotesize\bfseries\color{ink}（1） 承接点题}\par\vspace{1pt}
\begin{list}{}{\leftmargin=1.3em\labelwidth=0.95em\labelsep=0.35em\itemindent=0pt\topsep=0pt\partopsep=0pt\parsep=1.5pt\itemsep=0pt}
\item[{\color{accent}\scriptsize$\blacktriangleright$}]{\footnotesize You asked for my advice on {\color{accent}\bfseries 【主题】}. Having thought it over, I would suggest {\color{accent}\bfseries 【结论】}.}\par\vspace{0.5pt}{\footnotesize\color{gray}先给结论再解释，符合英语“先总后分”的行文习惯。}
\item[{\color{accent}\scriptsize$\blacktriangleright$}]{\footnotesize I am glad to hear that {\color{accent}\bfseries 【好消息】}, and I would like to offer a few suggestions on {\color{accent}\bfseries 【主题】}.}\par\vspace{0.5pt}{\footnotesize\color{gray}专用于“祝贺＋建议”合并题（2011），一句把两件事都点到。}
\end{list}
\vspace{2pt}{\footnotesize\bfseries\color{ink}（2） 建议 + 理由}\par\vspace{1pt}
\begin{list}{}{\leftmargin=1.3em\labelwidth=0.95em\labelsep=0.35em\itemindent=0pt\topsep=0pt\partopsep=0pt\parsep=1.5pt\itemsep=0pt}
\item[{\color{accent}\scriptsize$\blacktriangleright$}]{\footnotesize For one thing, it might be a good idea to {\color{accent}\bfseries 【建议一】}, since {\color{accent}\bfseries 【理由】}.}\par\vspace{0.5pt}{\footnotesize\color{gray}it might be a good idea to 比 you should 委婉得多。}
\item[{\color{accent}\scriptsize$\blacktriangleright$}]{\footnotesize For another, {\color{accent}\bfseries 【建议二】} would also help you {\color{accent}\bfseries 【效果】}.}\par\vspace{0.5pt}{\footnotesize\color{gray}第二条用 would also help，避免与 for one thing 重复句式。}
\end{list}
\vspace{2pt}{\footnotesize\bfseries\color{ink}（3） 收尾}\par\vspace{1pt}
\begin{list}{}{\leftmargin=1.3em\labelwidth=0.95em\labelsep=0.35em\itemindent=0pt\topsep=0pt\partopsep=0pt\parsep=1.5pt\itemsep=0pt}
\item[{\color{accent}\scriptsize$\blacktriangleright$}]{\footnotesize I hope you will find these suggestions helpful.}\par\vspace{0.5pt}{\footnotesize\color{gray}标准收尾，一句足够，不必再加三段客套。}
\item[{\color{accent}\scriptsize$\blacktriangleright$}]{\footnotesize Do let me know if you need further details. Best wishes for {\color{accent}\bfseries 【事项】}!}\par\vspace{0.5pt}{\footnotesize\color{gray}留一个开放的帮助姿态，比 I hope you are happy 具体。}
\end{list}
\vspace{1pt}{\footnotesize\color{gray}语气与落款：一封信里命令句（you must / you should）不要超过一句，建议信最忌命令语气。}\par
\endgroup
\vspace{4pt}\noindent{\color{accent}\bfseries 四、满分范文}\par\nopagebreak
{\color{rule}\rule{\linewidth}{0.6pt}}\par\nopagebreak
\vspace{1pt}
\vspace{2pt}\noindent\colorbox{soft}{\parbox{\dimexpr\linewidth-2\fboxsep\relax}{%
\vspace{2pt}
{\footnotesize\bfseries\color{accent}2023 真题话题：建议朋友选择看哪个展览}\hfill{\footnotesize\color{gray}正文 93 词}
\par\vspace{2pt}
\footnotesize Dear Mike,\par
\footnotesize You asked for my advice on which exhibition to visit this weekend. Having thought it over, I would suggest the art exhibition.\par
\footnotesize For one thing, it displays works by young Chinese painters, and they are far more original than the robot models at the other exhibition. For another, the gallery is only ten minutes' walk from your hotel, so you will not waste time on the way there. If you go on Sunday morning, you can also join a free guided tour.\par
\footnotesize I hope you will find these suggestions helpful. Enjoy your weekend!\par
\footnotesize Yours sincerely,\par
\footnotesize Li Ming\par
\par\vspace{2pt}{\footnotesize\color{gray}亲爱的迈克：你问我这个周末该去看哪个展览。我想了想，建议你去艺术展。一方面，它展出中国青年画家的作品，比另一个展览上的机器人模型有创意得多。另一方面，美术馆离你酒店只有十分钟步行路程，不会把时间浪费在路上。如果周日上午去，还能参加免费的导览。希望这些建议对你有用。周末愉快！你真诚的，李明}
\par\vspace{1pt}}}\par
\vspace{2pt}
\vspace{4pt}\noindent{\color{accent}\bfseries 五、雷区与得分点}\par\nopagebreak
{\color{rule}\rule{\linewidth}{0.6pt}}\par\nopagebreak
\vspace{1pt}
\vspace{1pt}\begingroup\setlength{\parindent}{0pt}
\begin{list}{}{\leftmargin=1.3em\labelwidth=0.95em\labelsep=0.35em\itemindent=0pt\topsep=0pt\partopsep=0pt\parsep=1.5pt\itemsep=0pt}
\item[{\color{warn}\scriptsize$\blacksquare$}]{\footnotesize 建议只有“多读书、多练习”这类空话，没有理由，内容单薄，落在 6 分档。}
\item[{\color{warn}\scriptsize$\blacksquare$}]{\footnotesize 两条建议其实是一条（都说“多练”），等于只完成一个要点。}
\item[{\color{warn}\scriptsize$\blacksquare$}]{\footnotesize 合并题漏答：2011 要求“祝贺＋建议”，只写建议就丢掉一半任务。}
\item[{\color{warn}\scriptsize$\blacksquare$}]{\footnotesize 得分点：先给结论（suggest the art exhibition）→ 两条带理由的建议 → 一条可执行的具体做法（周日上午有免费导览）。}
\end{list}
\endgroup
\clearpage
\label{g04}
\noindent\colorbox{accent}{%
\parbox{\dimexpr\linewidth-2\fboxsep\relax}{%
\color{white}\vspace{1pt}
{\LARGE\bfseries 04}\hspace{0.5em}{\Large 感谢信}\hfill{\large\itshape Letter of Thanks}
\\[1pt]{\footnotesize 真题 2010 · 2016}\par\vspace{1pt}}}
\par
\vspace{4pt}\noindent\colorbox{soft}{%
\parbox{\dimexpr\linewidth-2\fboxsep\relax}{\footnotesize{\color{accent}\bfseries 识别信号}\quad 题干出现 thank / grateful / 感谢……的帮助或接待，通常还要求“说明这次帮助带来的影响”或“表示回报”。}}
\par
\vspace{2pt}
\vspace{4pt}\noindent{\color{accent}\bfseries 一、任务要点}\par\nopagebreak
{\color{rule}\rule{\linewidth}{0.6pt}}\par\nopagebreak
\vspace{1pt}
\vspace{2pt}\noindent\begin{tabularx}{\dimexpr\linewidth-1.4pt\relax}{@{}>{\raggedright\arraybackslash}p{2.1cm}X@{}}
\toprule
\footnotesize{\color{accent}\bfseries 写作身份} & \footnotesize 受助者，回信给帮助过自己的人 \\
\midrule
\footnotesize{\color{accent}\bfseries 必须写到} & \footnotesize （1）为哪件事致谢（具体）　（2）一个细节 + 没有对方帮助会怎样　（3）再次致谢或表示回报 \\
\midrule
\footnotesize{\color{accent}\bfseries 一句话策略} & \footnotesize 感谢信的说服力全在“具体”：写出一件事、一个动作，比写十遍 thank you 有效 \\
\bottomrule
\end{tabularx}\par
\vspace{2pt}
\vspace{4pt}\noindent{\color{accent}\bfseries 二、结构骨架}\par\nopagebreak
{\color{rule}\rule{\linewidth}{0.6pt}}\par\nopagebreak
\vspace{1pt}
\vspace{2pt}\noindent\begingroup\footnotesize\renewcommand{\arraystretch}{1.15}
\begin{tabularx}{\dimexpr\linewidth-1.4pt\relax}{@{}>{\raggedright\arraybackslash}p{2.4cm}>{\raggedright\arraybackslash}p{3.4cm}>{\raggedright\arraybackslash}X@{}}
\toprule
{\color{accent}\bfseries 位置} & {\color{accent}\bfseries 功能} & {\color{accent}\bfseries 必须落到纸面的内容} \\
\midrule
称呼 & 按关系定 & Dear Ms. Brown, / Dear Jack, \\
第一段 & 点题 & 为哪件事致谢，一句话说清，不铺垫 \\
第二段 & 具体化 & 回忆一个细节 + 没有帮助会怎样 / 帮助带来的结果 \\
第三段 & 收尾 & 再次感谢 + 希望有机会回报 \\
落款 & 格式 & Yours sincerely, + Li Ming \\
\bottomrule
\end{tabularx}\endgroup\par
\vspace{2pt}
\vspace{4pt}\noindent{\color{accent}\bfseries 三、可替换骨架句}\par\nopagebreak
{\color{rule}\rule{\linewidth}{0.6pt}}\par\nopagebreak
\vspace{1pt}
\vspace{1pt}\begingroup\setlength{\parindent}{0pt}
\vspace{2pt}{\footnotesize\bfseries\color{ink}（1） 点题}\par\vspace{1pt}
\begin{list}{}{\leftmargin=1.3em\labelwidth=0.95em\labelsep=0.35em\itemindent=0pt\topsep=0pt\partopsep=0pt\parsep=1.5pt\itemsep=0pt}
\item[{\color{accent}\scriptsize$\blacktriangleright$}]{\footnotesize I am writing to express my sincere gratitude for {\color{accent}\bfseries 【帮助】}.}\par\vspace{0.5pt}{\footnotesize\color{gray}最稳的一句；gratitude 后接 for，不接 to。}
\item[{\color{accent}\scriptsize$\blacktriangleright$}]{\footnotesize I would like to thank you once again for {\color{accent}\bfseries 【帮助】} during {\color{accent}\bfseries 【时段】}.}\par\vspace{0.5pt}{\footnotesize\color{gray}用于回访、回信场景，once again 表明不是第一次致谢。}
\end{list}
\vspace{2pt}{\footnotesize\bfseries\color{ink}（2） 具体化}\par\vspace{1pt}
\begin{list}{}{\leftmargin=1.3em\labelwidth=0.95em\labelsep=0.35em\itemindent=0pt\topsep=0pt\partopsep=0pt\parsep=1.5pt\itemsep=0pt}
\item[{\color{accent}\scriptsize$\blacktriangleright$}]{\footnotesize What impressed me most was {\color{accent}\bfseries 【细节】}.}\par\vspace{0.5pt}{\footnotesize\color{gray}用 what 从句把感谢落到一件事上，避免空泛。}
\item[{\color{accent}\scriptsize$\blacktriangleright$}]{\footnotesize Without your help, I could hardly have {\color{accent}\bfseries 【结果】}.}\par\vspace{0.5pt}{\footnotesize\color{gray}虚拟语气，是本文体价值最高的句式，务必背下来。}
\end{list}
\vspace{2pt}{\footnotesize\bfseries\color{ink}（3） 收尾}\par\vspace{1pt}
\begin{list}{}{\leftmargin=1.3em\labelwidth=0.95em\labelsep=0.35em\itemindent=0pt\topsep=0pt\partopsep=0pt\parsep=1.5pt\itemsep=0pt}
\item[{\color{accent}\scriptsize$\blacktriangleright$}]{\footnotesize Once again, thank you for your kindness.}\par\vspace{0.5pt}{\footnotesize\color{gray}简洁收束，kindness 概括对方的好意。}
\item[{\color{accent}\scriptsize$\blacktriangleright$}]{\footnotesize I hope I can return the favour when you visit China.}\par\vspace{0.5pt}{\footnotesize\color{gray}给出回报的具体场景，比 I will repay you 自然。}
\end{list}
\vspace{1pt}{\footnotesize\color{gray}语气与落款：对师长、公务对象用前者；对朋友可用后者，语气更亲。}\par
\endgroup
\vspace{4pt}\noindent{\color{accent}\bfseries 四、满分范文}\par\nopagebreak
{\color{rule}\rule{\linewidth}{0.6pt}}\par\nopagebreak
\vspace{1pt}
\vspace{2pt}\noindent\colorbox{soft}{\parbox{\dimexpr\linewidth-2\fboxsep\relax}{%
\vspace{2pt}
{\footnotesize\bfseries\color{accent}2010 真题话题：感谢在美期间的接待}\hfill{\footnotesize\color{gray}正文 95 词}
\par\vspace{2pt}
\footnotesize Dear Ms. Brown,\par
\footnotesize I am writing to express my sincere gratitude for the warm reception you gave me during my stay in the United States.\par
\footnotesize What impressed me most was the trouble you took to show me around the campus and to introduce me to your colleagues. Thanks to your help, I soon felt at home in a strange city and finished my research on time. The books you lent me have also been of great use in my writing.\par
\footnotesize Once again, thank you for your kindness. I hope I can return the favour when you visit China.\par
\footnotesize Yours sincerely,\par
\footnotesize Li Ming\par
\par\vspace{2pt}{\footnotesize\color{gray}尊敬的布朗女士：我写信诚挚感谢您在我赴美期间给予的热情接待。最让我难忘的是您不辞辛劳带我参观校园，并把我介绍给您的同事。多亏您的帮助，我在陌生的城市很快安顿下来，并按时完成了研究。您借给我的书对我的写作也很有帮助。再次感谢您的好意。希望您来中国时我能回报这份情谊。您真诚的，李明}
\par\vspace{1pt}}}\par
\vspace{2pt}
\vspace{4pt}\noindent{\color{accent}\bfseries 五、雷区与得分点}\par\nopagebreak
{\color{rule}\rule{\linewidth}{0.6pt}}\par\nopagebreak
\vspace{1pt}
\vspace{1pt}\begingroup\setlength{\parindent}{0pt}
\begin{list}{}{\leftmargin=1.3em\labelwidth=0.95em\labelsep=0.35em\itemindent=0pt\topsep=0pt\partopsep=0pt\parsep=1.5pt\itemsep=0pt}
\item[{\color{warn}\scriptsize$\blacksquare$}]{\footnotesize 通篇 thank you very much 重复三次却不写具体事实，内容空，落在 5 分档。}
\item[{\color{warn}\scriptsize$\blacksquare$}]{\footnotesize 搭配错：gratitude to your help 应为 gratitude for；grateful to you for 才是对的。}
\item[{\color{warn}\scriptsize$\blacksquare$}]{\footnotesize 没交代“帮助的具体内容”，只写 your help，属于要点不完整。}
\item[{\color{warn}\scriptsize$\blacksquare$}]{\footnotesize 得分点：what 从句写细节 + without 虚拟语气写影响 + 结尾给回报的具体场景。}
\end{list}
\endgroup
\clearpage
\label{g05}
\noindent\colorbox{accent}{%
\parbox{\dimexpr\linewidth-2\fboxsep\relax}{%
\color{white}\vspace{1pt}
{\LARGE\bfseries 05}\hspace{0.5em}{\Large 祝贺信}\hfill{\large\itshape Letter of Congratulation}
\\[1pt]{\footnotesize 真题 2011}\par\vspace{1pt}}}
\par
\vspace{4pt}\noindent\colorbox{soft}{%
\parbox{\dimexpr\linewidth-2\fboxsep\relax}{\footnotesize{\color{accent}\bfseries 识别信号}\quad 题干出现 congratulate / 祝贺……考上、获奖、升职；英语二常与“建议”合并考，两个动作都要写。}}
\par
\vspace{2pt}
\vspace{4pt}\noindent{\color{accent}\bfseries 一、任务要点}\par\nopagebreak
{\color{rule}\rule{\linewidth}{0.6pt}}\par\nopagebreak
\vspace{1pt}
\vspace{2pt}\noindent\begin{tabularx}{\dimexpr\linewidth-1.4pt\relax}{@{}>{\raggedright\arraybackslash}p{2.1cm}X@{}}
\toprule
\footnotesize{\color{accent}\bfseries 写作身份} & \footnotesize 朋友、同学或同事，语气真诚但不过度吹捧 \\
\midrule
\footnotesize{\color{accent}\bfseries 必须写到} & \footnotesize （1）祝贺什么（具体成就）　（2）肯定对方的努力或品质（举一个具体记忆）　（3）祝愿或期待后续 \\
\midrule
\footnotesize{\color{accent}\bfseries 一句话策略} & \footnotesize 祝贺信最容易写成“恭喜恭喜”四句空话；把第二段写成“为什么这个成绩值得”，分数就稳了 \\
\bottomrule
\end{tabularx}\par
\vspace{2pt}
\vspace{4pt}\noindent{\color{accent}\bfseries 二、结构骨架}\par\nopagebreak
{\color{rule}\rule{\linewidth}{0.6pt}}\par\nopagebreak
\vspace{1pt}
\vspace{2pt}\noindent\begingroup\footnotesize\renewcommand{\arraystretch}{1.15}
\begin{tabularx}{\dimexpr\linewidth-1.4pt\relax}{@{}>{\raggedright\arraybackslash}p{2.4cm}>{\raggedright\arraybackslash}p{3.4cm}>{\raggedright\arraybackslash}X@{}}
\toprule
{\color{accent}\bfseries 位置} & {\color{accent}\bfseries 功能} & {\color{accent}\bfseries 必须落到纸面的内容} \\
\midrule
称呼 & 按关系定 & Dear Li Wei, / Dear Mr. Wang, \\
第一段 & 点题 & 听到喜讯 + Congratulations on…（on 不能换成 for） \\
第二段 & 肯定努力 & 这不是运气，是努力的结果，附一个具体记忆 \\
第三段 & 祝愿 & 祝未来更好 + 期待见面或来信 \\
落款 & 格式 & Yours sincerely, + Li Ming \\
\bottomrule
\end{tabularx}\endgroup\par
\vspace{2pt}
\vspace{4pt}\noindent{\color{accent}\bfseries 三、可替换骨架句}\par\nopagebreak
{\color{rule}\rule{\linewidth}{0.6pt}}\par\nopagebreak
\vspace{1pt}
\vspace{1pt}\begingroup\setlength{\parindent}{0pt}
\vspace{2pt}{\footnotesize\bfseries\color{ink}（1） 点题}\par\vspace{1pt}
\begin{list}{}{\leftmargin=1.3em\labelwidth=0.95em\labelsep=0.35em\itemindent=0pt\topsep=0pt\partopsep=0pt\parsep=1.5pt\itemsep=0pt}
\item[{\color{accent}\scriptsize$\blacktriangleright$}]{\footnotesize I was delighted to learn that {\color{accent}\bfseries 【喜讯】}. Congratulations on {\color{accent}\bfseries 【成就】}!}\par\vspace{0.5pt}{\footnotesize\color{gray}先陈述事实再祝贺，语序自然；congratulations 恒用复数。}
\item[{\color{accent}\scriptsize$\blacktriangleright$}]{\footnotesize Please accept my warmest congratulations on {\color{accent}\bfseries 【成就】}.}\par\vspace{0.5pt}{\footnotesize\color{gray}正式版，给同事、老师用。}
\end{list}
\vspace{2pt}{\footnotesize\bfseries\color{ink}（2） 肯定努力}\par\vspace{1pt}
\begin{list}{}{\leftmargin=1.3em\labelwidth=0.95em\labelsep=0.35em\itemindent=0pt\topsep=0pt\partopsep=0pt\parsep=1.5pt\itemsep=0pt}
\item[{\color{accent}\scriptsize$\blacktriangleright$}]{\footnotesize This success is by no means a matter of luck; it is the reward for {\color{accent}\bfseries 【努力】}.}\par\vspace{0.5pt}{\footnotesize\color{gray}by no means 加强语气，比 not 有力，也显出你真懂对方。}
\item[{\color{accent}\scriptsize$\blacktriangleright$}]{\footnotesize I still remember {\color{accent}\bfseries 【细节】}, and now it has paid off.}\par\vspace{0.5pt}{\footnotesize\color{gray}pay off 意为“得到回报”，是评价努力的标准搭配。}
\end{list}
\vspace{2pt}{\footnotesize\bfseries\color{ink}（3） 祝愿}\par\vspace{1pt}
\begin{list}{}{\leftmargin=1.3em\labelwidth=0.95em\labelsep=0.35em\itemindent=0pt\topsep=0pt\partopsep=0pt\parsep=1.5pt\itemsep=0pt}
\item[{\color{accent}\scriptsize$\blacktriangleright$}]{\footnotesize I wish you even greater progress in {\color{accent}\bfseries 【领域】}.}\par\vspace{0.5pt}{\footnotesize\color{gray}祝愿用 wish + 名词，不写 hope you will be better。}
\item[{\color{accent}\scriptsize$\blacktriangleright$}]{\footnotesize I am sure you will go further in the years ahead.}\par\vspace{0.5pt}{\footnotesize\color{gray}go further 比 be successful 具体，着眼未来。}
\end{list}
\vspace{1pt}{\footnotesize\color{gray}语气与落款：两个短句可以用在私人信第二段开头，但公务信里不宜连续口语化。}\par
\endgroup
\vspace{4pt}\noindent{\color{accent}\bfseries 四、满分范文}\par\nopagebreak
{\color{rule}\rule{\linewidth}{0.6pt}}\par\nopagebreak
\vspace{1pt}
\vspace{2pt}\noindent\colorbox{soft}{\parbox{\dimexpr\linewidth-2\fboxsep\relax}{%
\vspace{2pt}
{\footnotesize\bfseries\color{accent}2011 真题话题：祝贺朋友考上大学}\hfill{\footnotesize\color{gray}正文 92 词}
\par\vspace{2pt}
\footnotesize Dear Li Wei,\par
\footnotesize I was delighted to learn that you have been admitted to Peking University. Congratulations on this wonderful achievement!\par
\footnotesize Having worked side by side with you for three years, I know how much effort you put into your study. This success is by no means a matter of luck; it is the reward for the hours you spent in the library while others were playing games. Your parents must be very proud of you.\par
\footnotesize I wish you even greater progress in the years ahead. Do write and tell me about your new life.\par
\footnotesize Yours sincerely,\par
\footnotesize Li Ming\par
\par\vspace{2pt}{\footnotesize\color{gray}亲爱的李伟：得知你被北京大学录取，我非常高兴。祝贺你取得这样出色的成绩！与你并肩学习三年，我知道你付出了多少努力。这次成功绝不是运气，而是你在别人打游戏时泡在图书馆的那些时光换来的回报。你父母一定为你感到骄傲。祝你在未来几年取得更大进步。一定写信告诉我你的新生活。你真诚的，李明}
\par\vspace{1pt}}}\par
\vspace{2pt}
\vspace{4pt}\noindent{\color{accent}\bfseries 五、雷区与得分点}\par\nopagebreak
{\color{rule}\rule{\linewidth}{0.6pt}}\par\nopagebreak
\vspace{1pt}
\vspace{1pt}\begingroup\setlength{\parindent}{0pt}
\begin{list}{}{\leftmargin=1.3em\labelwidth=0.95em\labelsep=0.35em\itemindent=0pt\topsep=0pt\partopsep=0pt\parsep=1.5pt\itemsep=0pt}
\item[{\color{warn}\scriptsize$\blacksquare$}]{\footnotesize 只写 Congratulations! 一句，没有“为什么值得祝贺”，内容不足，5—6 分。}
\item[{\color{warn}\scriptsize$\blacksquare$}]{\footnotesize 搭配错：Congratulations for your success 应为 on your success。}
\item[{\color{warn}\scriptsize$\blacksquare$}]{\footnotesize 2011 是合并题，只写祝贺忘了建议，或只写建议忘了祝贺，都会丢要点。}
\item[{\color{warn}\scriptsize$\blacksquare$}]{\footnotesize 得分点：具体成就 + 一个共同记忆 + by no means 的肯定 + 面向未来的祝愿。}
\end{list}
\endgroup
\clearpage
\label{g06}
\noindent\colorbox{accent}{%
\parbox{\dimexpr\linewidth-2\fboxsep\relax}{%
\color{white}\vspace{1pt}
{\LARGE\bfseries 06}\hspace{0.5em}{\Large 道歉信}\hfill{\large\itshape Letter of Apology}
\\[1pt]{\footnotesize 真题 2018}\par\vspace{1pt}}}
\par
\vspace{4pt}\noindent\colorbox{soft}{%
\parbox{\dimexpr\linewidth-2\fboxsep\relax}{\footnotesize{\color{accent}\bfseries 识别信号}\quad 题干出现 apologize / 为……道歉，通常要求“说明原因”并“提出补救办法”；只写对不起不给补救会丢要点。}}
\par
\vspace{2pt}
\vspace{4pt}\noindent{\color{accent}\bfseries 一、任务要点}\par\nopagebreak
{\color{rule}\rule{\linewidth}{0.6pt}}\par\nopagebreak
\vspace{1pt}
\vspace{2pt}\noindent\begin{tabularx}{\dimexpr\linewidth-1.4pt\relax}{@{}>{\raggedright\arraybackslash}p{2.1cm}X@{}}
\toprule
\footnotesize{\color{accent}\bfseries 写作身份} & \footnotesize 失约、失误的一方，态度要主动担责 \\
\midrule
\footnotesize{\color{accent}\bfseries 必须写到} & \footnotesize （1）为什么事道歉　（2）客观原因（一句，不推卸）　（3）补救方案（改期、替代、赔偿）　（4）再次致歉 \\
\midrule
\footnotesize{\color{accent}\bfseries 一句话策略} & \footnotesize 道歉信的结构是“认错—解释—补救—再致歉”，解释只占一句，补救必须具体到时间 \\
\bottomrule
\end{tabularx}\par
\vspace{2pt}
\vspace{4pt}\noindent{\color{accent}\bfseries 二、结构骨架}\par\nopagebreak
{\color{rule}\rule{\linewidth}{0.6pt}}\par\nopagebreak
\vspace{1pt}
\vspace{2pt}\noindent\begingroup\footnotesize\renewcommand{\arraystretch}{1.15}
\begin{tabularx}{\dimexpr\linewidth-1.4pt\relax}{@{}>{\raggedright\arraybackslash}p{2.4cm}>{\raggedright\arraybackslash}p{3.4cm}>{\raggedright\arraybackslash}X@{}}
\toprule
{\color{accent}\bfseries 位置} & {\color{accent}\bfseries 功能} & {\color{accent}\bfseries 必须落到纸面的内容} \\
\midrule
称呼 & 按关系定 & Dear Professor Wang, / Dear Jack, \\
第一段 & 点题 & apologize for + 事项，一句到位，不先讲一堆理由 \\
第二段 & 原因 + 补救 & 原因一句 + 具体补救（改期到哪天、用什么替代） \\
第三段 & 收尾 & 再次致歉 + 感谢理解 / 期待弥补 \\
落款 & 格式 & Yours sincerely, + Li Ming \\
\bottomrule
\end{tabularx}\endgroup\par
\vspace{2pt}
\vspace{4pt}\noindent{\color{accent}\bfseries 三、可替换骨架句}\par\nopagebreak
{\color{rule}\rule{\linewidth}{0.6pt}}\par\nopagebreak
\vspace{1pt}
\vspace{1pt}\begingroup\setlength{\parindent}{0pt}
\vspace{2pt}{\footnotesize\bfseries\color{ink}（1） 点题}\par\vspace{1pt}
\begin{list}{}{\leftmargin=1.3em\labelwidth=0.95em\labelsep=0.35em\itemindent=0pt\topsep=0pt\partopsep=0pt\parsep=1.5pt\itemsep=0pt}
\item[{\color{accent}\scriptsize$\blacktriangleright$}]{\footnotesize I am writing to apologize for {\color{accent}\bfseries 【事项】}. Please accept my sincere apologies.}\par\vspace{0.5pt}{\footnotesize\color{gray}两句一组：先说明事由，再正式致歉。}
\item[{\color{accent}\scriptsize$\blacktriangleright$}]{\footnotesize I feel terribly sorry that I have to {\color{accent}\bfseries 【取消/推迟】} {\color{accent}\bfseries 【事项】}.}\par\vspace{0.5pt}{\footnotesize\color{gray}需要表达“不得已”时用 have to，暗示并非本意。}
\end{list}
\vspace{2pt}{\footnotesize\bfseries\color{ink}（2） 原因 + 补救}\par\vspace{1pt}
\begin{list}{}{\leftmargin=1.3em\labelwidth=0.95em\labelsep=0.35em\itemindent=0pt\topsep=0pt\partopsep=0pt\parsep=1.5pt\itemsep=0pt}
\item[{\color{accent}\scriptsize$\blacktriangleright$}]{\footnotesize The reason is that {\color{accent}\bfseries 【客观原因】}, which I cannot put off.}\par\vspace{0.5pt}{\footnotesize\color{gray}先说原因、再补一句“推不掉”，比堆三个理由真诚。}
\item[{\color{accent}\scriptsize$\blacktriangleright$}]{\footnotesize To make up for it, I would like to {\color{accent}\bfseries 【补救】} instead.}\par\vspace{0.5pt}{\footnotesize\color{gray}make up for 是“补救”的标准说法，后接名词或动名词。}
\end{list}
\vspace{2pt}{\footnotesize\bfseries\color{ink}（3） 收尾}\par\vspace{1pt}
\begin{list}{}{\leftmargin=1.3em\labelwidth=0.95em\labelsep=0.35em\itemindent=0pt\topsep=0pt\partopsep=0pt\parsep=1.5pt\itemsep=0pt}
\item[{\color{accent}\scriptsize$\blacktriangleright$}]{\footnotesize I am sorry for any inconvenience I have caused.}\par\vspace{0.5pt}{\footnotesize\color{gray}固定客套句，公务道歉信必写，一句即可。}
\item[{\color{accent}\scriptsize$\blacktriangleright$}]{\footnotesize Thank you for your understanding, and I hope we can meet soon.}\par\vspace{0.5pt}{\footnotesize\color{gray}道歉信不能停在“对不起”，结尾要给下一步。}
\end{list}
\vspace{1pt}{\footnotesize\color{gray}语气与落款：do 加强语气可用于私人信；公务信与其加重语气，不如把补救写具体。}\par
\endgroup
\vspace{4pt}\noindent{\color{accent}\bfseries 四、满分范文}\par\nopagebreak
{\color{rule}\rule{\linewidth}{0.6pt}}\par\nopagebreak
\vspace{1pt}
\vspace{2pt}\noindent\colorbox{soft}{\parbox{\dimexpr\linewidth-2\fboxsep\relax}{%
\vspace{2pt}
{\footnotesize\bfseries\color{accent}2018 真题话题：为临时取消拜访道歉并建议改期}\hfill{\footnotesize\color{gray}正文 92 词}
\par\vspace{2pt}
\footnotesize Dear Professor Wang,\par
\footnotesize I am writing to apologize for cancelling my visit to your office this Friday. Please accept my sincere apologies.\par
\footnotesize The reason is that I have to attend an emergency meeting in Shanghai, which I cannot put off. To make up for it, I would like to come next Monday morning instead, and I will bring the materials we planned to discuss. Would ten o'clock suit you? I will call your office on Friday to confirm the new time.\par
\footnotesize I am sorry for any inconvenience I have caused. Thank you for your understanding.\par
\footnotesize Yours sincerely,\par
\footnotesize Li Ming\par
\par\vspace{2pt}{\footnotesize\color{gray}尊敬的王教授：我写信为取消本周五到您办公室的拜访道歉，请接受我诚挚的歉意。原因是我要去上海参加一个无法推迟的紧急会议。为弥补，我想改为下周一上午前来，并带上我们原计划讨论的材料。十点对您方便吗？周五我会打电话到您办公室确认新的时间。给您的任何不便，我深表歉意。感谢您的理解。您真诚的，李明}
\par\vspace{1pt}}}\par
\vspace{2pt}
\vspace{4pt}\noindent{\color{accent}\bfseries 五、雷区与得分点}\par\nopagebreak
{\color{rule}\rule{\linewidth}{0.6pt}}\par\nopagebreak
\vspace{1pt}
\vspace{1pt}\begingroup\setlength{\parindent}{0pt}
\begin{list}{}{\leftmargin=1.3em\labelwidth=0.95em\labelsep=0.35em\itemindent=0pt\topsep=0pt\partopsep=0pt\parsep=1.5pt\itemsep=0pt}
\item[{\color{warn}\scriptsize$\blacksquare$}]{\footnotesize 原因写成三条（堵车、天气、家人住院），越解释越像借口；一句足够。}
\item[{\color{warn}\scriptsize$\blacksquare$}]{\footnotesize 只道歉不给补救，任务里的“提出解决办法”缺失，顶多 5—6 分。}
\item[{\color{warn}\scriptsize$\blacksquare$}]{\footnotesize 用 I am sorry to… 代替 I apologize for…：前者是“遗憾”，后者才是“道歉”。}
\item[{\color{warn}\scriptsize$\blacksquare$}]{\footnotesize 得分点：apologize for 点题 + 一句客观原因 + 具体改期（next Monday, ten o'clock）+ 再次致歉收尾。}
\end{list}
\endgroup
\clearpage
\label{g07}
\noindent\colorbox{accent}{%
\parbox{\dimexpr\linewidth-2\fboxsep\relax}{%
\color{white}\vspace{1pt}
{\LARGE\bfseries 07}\hspace{0.5em}{\Large 投诉信}\hfill{\large\itshape Letter of Complaint}
\\[1pt]{\footnotesize 真题 2012}\par\vspace{1pt}}}
\par
\vspace{4pt}\noindent\colorbox{soft}{%
\parbox{\dimexpr\linewidth-2\fboxsep\relax}{\footnotesize{\color{accent}\bfseries 识别信号}\quad 题干出现 complain / 投诉商品或服务，通常要求“说明问题”并“提出解决方式”，部分题目还要求给出期限。}}
\par
\vspace{2pt}
\vspace{4pt}\noindent{\color{accent}\bfseries 一、任务要点}\par\nopagebreak
{\color{rule}\rule{\linewidth}{0.6pt}}\par\nopagebreak
\vspace{1pt}
\vspace{2pt}\noindent\begin{tabularx}{\dimexpr\linewidth-1.4pt\relax}{@{}>{\raggedright\arraybackslash}p{2.1cm}X@{}}
\toprule
\footnotesize{\color{accent}\bfseries 写作身份} & \footnotesize 消费者或使用者，语气要有理有据，不是发泄 \\
\midrule
\footnotesize{\color{accent}\bfseries 必须写到} & \footnotesize （1）投诉对象（何时、从何渠道购买）　（2）具体问题 + 已做过的尝试　（3）明确要求（退款/更换/赔偿）＋期限 \\
\midrule
\footnotesize{\color{accent}\bfseries 一句话策略} & \footnotesize 事实三句、要求一句。写清“我试过什么”，投诉才站得住脚 \\
\bottomrule
\end{tabularx}\par
\vspace{2pt}
\vspace{4pt}\noindent{\color{accent}\bfseries 二、结构骨架}\par\nopagebreak
{\color{rule}\rule{\linewidth}{0.6pt}}\par\nopagebreak
\vspace{1pt}
\vspace{2pt}\noindent\begingroup\footnotesize\renewcommand{\arraystretch}{1.15}
\begin{tabularx}{\dimexpr\linewidth-1.4pt\relax}{@{}>{\raggedright\arraybackslash}p{2.4cm}>{\raggedright\arraybackslash}p{3.4cm}>{\raggedright\arraybackslash}X@{}}
\toprule
{\color{accent}\bfseries 位置} & {\color{accent}\bfseries 功能} & {\color{accent}\bfseries 必须落到纸面的内容} \\
\midrule
称呼 & 不知收信人姓名时 & Dear Sir or Madam, \\
第一段 & 点题 & 投诉什么 + 购买时间与渠道，一句交代完 \\
第二段 & 问题 & 具体故障或失误 + 已做的尝试（换电池、按说明书操作） \\
第三段 & 要求 & 明确要求（replace / refund）＋期限＋否则将…… \\
落款 & 格式 & Yours sincerely, + Li Ming（若不知对方性别，可用 Yours faithfully,） \\
\bottomrule
\end{tabularx}\endgroup\par
\vspace{2pt}
\vspace{4pt}\noindent{\color{accent}\bfseries 三、可替换骨架句}\par\nopagebreak
{\color{rule}\rule{\linewidth}{0.6pt}}\par\nopagebreak
\vspace{1pt}
\vspace{1pt}\begingroup\setlength{\parindent}{0pt}
\vspace{2pt}{\footnotesize\bfseries\color{ink}（1） 点题}\par\vspace{1pt}
\begin{list}{}{\leftmargin=1.3em\labelwidth=0.95em\labelsep=0.35em\itemindent=0pt\topsep=0pt\partopsep=0pt\parsep=1.5pt\itemsep=0pt}
\item[{\color{accent}\scriptsize$\blacktriangleright$}]{\footnotesize I am writing to complain about {\color{accent}\bfseries 【商品/服务】}, which I ordered from {\color{accent}\bfseries 【渠道】} on {\color{accent}\bfseries 【日期】}.}\par\vspace{0.5pt}{\footnotesize\color{gray}一句交代对象、渠道、时间，信息密度最高。}
\item[{\color{accent}\scriptsize$\blacktriangleright$}]{\footnotesize I regret to say that I am far from satisfied with {\color{accent}\bfseries 【商品/服务】}.}\par\vspace{0.5pt}{\footnotesize\color{gray}far from satisfied 比 not satisfied 委婉但明确。}
\end{list}
\vspace{2pt}{\footnotesize\bfseries\color{ink}（2） 问题}\par\vspace{1pt}
\begin{list}{}{\leftmargin=1.3em\labelwidth=0.95em\labelsep=0.35em\itemindent=0pt\topsep=0pt\partopsep=0pt\parsep=1.5pt\itemsep=0pt}
\item[{\color{accent}\scriptsize$\blacktriangleright$}]{\footnotesize The problem is that {\color{accent}\bfseries 【具体故障】}.}\par\vspace{0.5pt}{\footnotesize\color{gray}The problem is that 后接完整句，不要接短语。}
\item[{\color{accent}\scriptsize$\blacktriangleright$}]{\footnotesize I have {\color{accent}\bfseries 【尝试】}, but the trouble remains.}\par\vspace{0.5pt}{\footnotesize\color{gray}写“已尝试”是投诉信的关键细节，说明不是自己不会用。}
\end{list}
\vspace{2pt}{\footnotesize\bfseries\color{ink}（3） 要求}\par\vspace{1pt}
\begin{list}{}{\leftmargin=1.3em\labelwidth=0.95em\labelsep=0.35em\itemindent=0pt\topsep=0pt\partopsep=0pt\parsep=1.5pt\itemsep=0pt}
\item[{\color{accent}\scriptsize$\blacktriangleright$}]{\footnotesize I would like you to replace it with a new one, or give me a full refund.}\par\vspace{0.5pt}{\footnotesize\color{gray}给两个可选方案，对方更容易接受，也显得理性。}
\item[{\color{accent}\scriptsize$\blacktriangleright$}]{\footnotesize I would be grateful if you could deal with this matter within {\color{accent}\bfseries 【期限】}.}\par\vspace{0.5pt}{\footnotesize\color{gray}期限是投诉信的“牙齿”，十天、两周都可以，但必须写。}
\end{list}
\vspace{1pt}{\footnotesize\color{gray}语气与落款：知道对方姓名用前者，Dear Sir or Madam 开头时用后者，二者不可混用。}\par
\endgroup
\vspace{4pt}\noindent{\color{accent}\bfseries 四、满分范文}\par\nopagebreak
{\color{rule}\rule{\linewidth}{0.6pt}}\par\nopagebreak
\vspace{1pt}
\vspace{2pt}\noindent\colorbox{soft}{\parbox{\dimexpr\linewidth-2\fboxsep\relax}{%
\vspace{2pt}
{\footnotesize\bfseries\color{accent}2012 真题话题：投诉网购电子词典}\hfill{\footnotesize\color{gray}正文 101 词}
\par\vspace{2pt}
\footnotesize Dear Sir or Madam,\par
\footnotesize I am writing to complain about the electronic dictionary I ordered from your online store on 10 May.\par
\footnotesize The problem is that the screen goes black whenever I turn the device on, and the battery runs out within an hour. I have changed the batteries twice and followed the instructions in the manual, but the trouble remains. As the dictionary is still under guarantee, I would like you to replace it with a new one or give me a full refund.\par
\footnotesize I would be grateful if you could deal with this matter within ten days. I look forward to your early reply.\par
\footnotesize Yours faithfully,\par
\footnotesize Li Ming\par
\par\vspace{2pt}{\footnotesize\color{gray}尊敬的先生/女士：我写信投诉 5 月 10 日从贵网店订购的电子词典。问题是每次开机屏幕都会变黑，电池一小时就用完。我已换过两次电池，并按说明书操作，但故障依旧。由于词典仍在保修期内，我希望你们更换新品或全额退款。如能在十天内处理此事，我将不胜感激。期待早日答复。您忠实的，李明}
\par\vspace{1pt}}}\par
\vspace{2pt}
\vspace{4pt}\noindent{\color{accent}\bfseries 五、雷区与得分点}\par\nopagebreak
{\color{rule}\rule{\linewidth}{0.6pt}}\par\nopagebreak
\vspace{1pt}
\vspace{1pt}\begingroup\setlength{\parindent}{0pt}
\begin{list}{}{\leftmargin=1.3em\labelwidth=0.95em\labelsep=0.35em\itemindent=0pt\topsep=0pt\partopsep=0pt\parsep=1.5pt\itemsep=0pt}
\item[{\color{warn}\scriptsize$\blacksquare$}]{\footnotesize 通篇情绪词（terrible, angry, cheated），不写事实，语言分与内容分双失。}
\item[{\color{warn}\scriptsize$\blacksquare$}]{\footnotesize 不写购买时间或渠道，对方无法查证，要点不全。}
\item[{\color{warn}\scriptsize$\blacksquare$}]{\footnotesize 一次性索要退款＋赔偿＋道歉＋解释，不现实，显得在背模板。}
\item[{\color{warn}\scriptsize$\blacksquare$}]{\footnotesize 得分点：The problem is that… 写具体故障 + I have…but… 写尝试 + 明确要求与期限。}
\end{list}
\endgroup
\clearpage
\label{g08}
\noindent\colorbox{accent}{%
\parbox{\dimexpr\linewidth-2\fboxsep\relax}{%
\color{white}\vspace{1pt}
{\LARGE\bfseries 08}\hspace{0.5em}{\Large 咨询信}\hfill{\large\itshape Letter of Inquiry}
\\[1pt]{\footnotesize 真题 2014}\par\vspace{1pt}}}
\par
\vspace{4pt}\noindent\colorbox{soft}{%
\parbox{\dimexpr\linewidth-2\fboxsep\relax}{\footnotesize{\color{accent}\bfseries 识别信号}\quad 题干出现 information / ask about / 询问……详情；2014 是“告知自己情况 + 咨询对方意见”的合并题。}}
\par
\vspace{2pt}
\vspace{4pt}\noindent{\color{accent}\bfseries 一、任务要点}\par\nopagebreak
{\color{rule}\rule{\linewidth}{0.6pt}}\par\nopagebreak
\vspace{1pt}
\vspace{2pt}\noindent\begin{tabularx}{\dimexpr\linewidth-1.4pt\relax}{@{}>{\raggedright\arraybackslash}p{2.1cm}X@{}}
\toprule
\footnotesize{\color{accent}\bfseries 写作身份} & \footnotesize 潜在参加者、申请者或新生；语气客气，问题具体 \\
\midrule
\footnotesize{\color{accent}\bfseries 必须写到} & \footnotesize （1）从何渠道了解到该项目　（2）2—3 个具体问题　（3）礼貌收尾（期待回复/可提供更多信息） \\
\midrule
\footnotesize{\color{accent}\bfseries 一句话策略} & \footnotesize “详情”不是问题，具体问题才是。时间、费用、条件、安排四类里挑两个问 \\
\bottomrule
\end{tabularx}\par
\vspace{2pt}
\vspace{4pt}\noindent{\color{accent}\bfseries 二、结构骨架}\par\nopagebreak
{\color{rule}\rule{\linewidth}{0.6pt}}\par\nopagebreak
\vspace{1pt}
\vspace{2pt}\noindent\begingroup\footnotesize\renewcommand{\arraystretch}{1.15}
\begin{tabularx}{\dimexpr\linewidth-1.4pt\relax}{@{}>{\raggedright\arraybackslash}p{2.4cm}>{\raggedright\arraybackslash}p{3.4cm}>{\raggedright\arraybackslash}X@{}}
\toprule
{\color{accent}\bfseries 位置} & {\color{accent}\bfseries 功能} & {\color{accent}\bfseries 必须落到纸面的内容} \\
\midrule
称呼 & 不知姓名时 & Dear Sir or Madam, \\
第一段 & 点题 & 从哪知道的 + 想咨询什么（项目名照抄题干） \\
第二段 & 提问 & 2—3 个具体问题：起止时间、费用、住宿、条件、报名方式 \\
第三段 & 收尾 & 期待回复 + 表示感谢，可留联系方式 \\
落款 & 格式 & Yours sincerely, + Li Ming \\
\bottomrule
\end{tabularx}\endgroup\par
\vspace{2pt}
\vspace{4pt}\noindent{\color{accent}\bfseries 三、可替换骨架句}\par\nopagebreak
{\color{rule}\rule{\linewidth}{0.6pt}}\par\nopagebreak
\vspace{1pt}
\vspace{1pt}\begingroup\setlength{\parindent}{0pt}
\vspace{2pt}{\footnotesize\bfseries\color{ink}（1） 点题}\par\vspace{1pt}
\begin{list}{}{\leftmargin=1.3em\labelwidth=0.95em\labelsep=0.35em\itemindent=0pt\topsep=0pt\partopsep=0pt\parsep=1.5pt\itemsep=0pt}
\item[{\color{accent}\scriptsize$\blacktriangleright$}]{\footnotesize I am writing to ask for some information about {\color{accent}\bfseries 【项目】}, which I learned about {\color{accent}\bfseries 【渠道】}.}\par\vspace{0.5pt}{\footnotesize\color{gray}交代信息来源，信才可信；which 从句省掉一个句子。}
\item[{\color{accent}\scriptsize$\blacktriangleright$}]{\footnotesize I am interested in {\color{accent}\bfseries 【项目】} and would like to know more about it.}\par\vspace{0.5pt}{\footnotesize\color{gray}先说兴趣、再提问，语气不生硬。}
\end{list}
\vspace{2pt}{\footnotesize\bfseries\color{ink}（2） 提问}\par\vspace{1pt}
\begin{list}{}{\leftmargin=1.3em\labelwidth=0.95em\labelsep=0.35em\itemindent=0pt\topsep=0pt\partopsep=0pt\parsep=1.5pt\itemsep=0pt}
\item[{\color{accent}\scriptsize$\blacktriangleright$}]{\footnotesize First, I would like to know when {\color{accent}\bfseries 【项目】} begins and how long it lasts.}\par\vspace{0.5pt}{\footnotesize\color{gray}问时间不能只写 when，要带“持续多久”，信息才具体。}
\item[{\color{accent}\scriptsize$\blacktriangleright$}]{\footnotesize Besides, could you tell me whether {\color{accent}\bfseries 【条件】} and how much {\color{accent}\bfseries 【费用】}?}\par\vspace{0.5pt}{\footnotesize\color{gray}whether 与 how much 并列，一次问两件事，句子不显啰嗦。}
\end{list}
\vspace{2pt}{\footnotesize\bfseries\color{ink}（3） 收尾}\par\vspace{1pt}
\begin{list}{}{\leftmargin=1.3em\labelwidth=0.95em\labelsep=0.35em\itemindent=0pt\topsep=0pt\partopsep=0pt\parsep=1.5pt\itemsep=0pt}
\item[{\color{accent}\scriptsize$\blacktriangleright$}]{\footnotesize I would appreciate it if you could reply at your earliest convenience.}\par\vspace{0.5pt}{\footnotesize\color{gray}标准公务收尾，at your earliest convenience 比 as soon as possible 客气。}
\item[{\color{accent}\scriptsize$\blacktriangleright$}]{\footnotesize Please let me know if you need any further information from me.}\par\vspace{0.5pt}{\footnotesize\color{gray}反向提供信息，显得配合、有诚意。}
\end{list}
\vspace{1pt}{\footnotesize\color{gray}语气与落款：问题较敏感（费用、名额）时用这句，比直接问更得体。}\par
\endgroup
\vspace{4pt}\noindent{\color{accent}\bfseries 四、满分范文}\par\nopagebreak
{\color{rule}\rule{\linewidth}{0.6pt}}\par\nopagebreak
\vspace{1pt}
\vspace{2pt}\noindent\colorbox{soft}{\parbox{\dimexpr\linewidth-2\fboxsep\relax}{%
\vspace{2pt}
{\footnotesize\bfseries\color{accent}2014 真题话题改编：咨询夏令营安排}\hfill{\footnotesize\color{gray}正文 91 词}
\par\vspace{2pt}
\footnotesize Dear Sir or Madam,\par
\footnotesize I am writing to ask for some information about the summer camp advertised on your website.\par
\footnotesize First, I would like to know when the camp begins and how long it lasts, as I have to arrange my travel. Besides, could you tell me whether accommodation is provided and how much it costs? I am also wondering if volunteers are needed, since I have experience in organizing campus activities.\par
\footnotesize I would appreciate it if you could reply at your earliest convenience. My telephone number is 12345678 in case you need to reach me.\par
\footnotesize Yours sincerely,\par
\footnotesize Li Ming\par
\par\vspace{2pt}{\footnotesize\color{gray}尊敬的先生/女士：我写信咨询贵网站宣传的夏令营的相关信息。首先，我想知道夏令营何时开始、持续多久，因为我需要安排行程。此外，能否告知是否提供住宿以及费用多少？我还想了解是否需要志愿者，因为我有组织校园活动的经验。如能尽早回复，我将不胜感激。如需联系我，电话是 12345678。您真诚的，李明}
\par\vspace{1pt}}}\par
\vspace{2pt}
\vspace{4pt}\noindent{\color{accent}\bfseries 五、雷区与得分点}\par\nopagebreak
{\color{rule}\rule{\linewidth}{0.6pt}}\par\nopagebreak
\vspace{1pt}
\vspace{1pt}\begingroup\setlength{\parindent}{0pt}
\begin{list}{}{\leftmargin=1.3em\labelwidth=0.95em\labelsep=0.35em\itemindent=0pt\topsep=0pt\partopsep=0pt\parsep=1.5pt\itemsep=0pt}
\item[{\color{warn}\scriptsize$\blacksquare$}]{\footnotesize 只问“详情”不列具体问题，对方无从回答，等于没完成任务。}
\item[{\color{warn}\scriptsize$\blacksquare$}]{\footnotesize 一封信用五个问号，显得没有重点；2—3 个足够。}
\item[{\color{warn}\scriptsize$\blacksquare$}]{\footnotesize 2014 要求先“告知自己的生活习惯”，漏掉告知部分，合并题失分。}
\item[{\color{warn}\scriptsize$\blacksquare$}]{\footnotesize 得分点：信息来源 + 两个具体问题（时间、住宿与费用）+ 一个补充问题 + 客气收尾。}
\end{list}
\endgroup
\clearpage
\label{g09}
\noindent\colorbox{accent}{%
\parbox{\dimexpr\linewidth-2\fboxsep\relax}{%
\color{white}\vspace{1pt}
{\LARGE\bfseries 09}\hspace{0.5em}{\Large 通知}\hfill{\large\itshape Notice}
\\[1pt]{\footnotesize 真题 2013 · 2015 · 2020}\par\vspace{1pt}}}
\par
\vspace{4pt}\noindent\colorbox{soft}{%
\parbox{\dimexpr\linewidth-2\fboxsep\relax}{\footnotesize{\color{accent}\bfseries 识别信号}\quad 题干出现 Notice / 通知……参加……，要求写一则通知；无称呼，必须有标题与落款。}}
\par
\vspace{2pt}
\vspace{4pt}\noindent{\color{accent}\bfseries 一、任务要点}\par\nopagebreak
{\color{rule}\rule{\linewidth}{0.6pt}}\par\nopagebreak
\vspace{1pt}
\vspace{2pt}\noindent\begin{tabularx}{\dimexpr\linewidth-1.4pt\relax}{@{}>{\raggedright\arraybackslash}p{2.1cm}X@{}}
\toprule
\footnotesize{\color{accent}\bfseries 写作身份} & \footnotesize 学生会、社团或组委会；用集体名义发布，不用 I \\
\midrule
\footnotesize{\color{accent}\bfseries 必须写到} & \footnotesize （1）活动目的　（2）时间、地点　（3）内容与参加方式　（4）报名截止或联系方式 \\
\midrule
\footnotesize{\color{accent}\bfseries 一句话策略} & \footnotesize 通知的第一句就要出现“为……目的 + 何时 + 何活动”，读者扫一眼就知道要不要来 \\
\bottomrule
\end{tabularx}\par
\vspace{2pt}
\vspace{4pt}\noindent{\color{accent}\bfseries 二、结构骨架}\par\nopagebreak
{\color{rule}\rule{\linewidth}{0.6pt}}\par\nopagebreak
\vspace{1pt}
\vspace{2pt}\noindent\begingroup\footnotesize\renewcommand{\arraystretch}{1.15}
\begin{tabularx}{\dimexpr\linewidth-1.4pt\relax}{@{}>{\raggedright\arraybackslash}p{2.4cm}>{\raggedright\arraybackslash}p{3.4cm}>{\raggedright\arraybackslash}X@{}}
\toprule
{\color{accent}\bfseries 位置} & {\color{accent}\bfseries 功能} & {\color{accent}\bfseries 必须落到纸面的内容} \\
\midrule
标题 & 居中单独一行 & Notice（首字母大写，不加句号，不必写日期） \\
第一段 & 目的 + 活动 & In order to…, the Student Union is to hold… on… \\
第二段 & 具体安排 & 起止时间、地点、活动内容、参加方式或报名要求 \\
第三段 & 号召 / 联系 & Everyone is welcome + 联系方式 \\
落款 & 右下角两行 & The Student Union / Oct. 20, 2025 \\
\bottomrule
\end{tabularx}\endgroup\par
\vspace{2pt}
\vspace{4pt}\noindent{\color{accent}\bfseries 三、可替换骨架句}\par\nopagebreak
{\color{rule}\rule{\linewidth}{0.6pt}}\par\nopagebreak
\vspace{1pt}
\vspace{1pt}\begingroup\setlength{\parindent}{0pt}
\vspace{2pt}{\footnotesize\bfseries\color{ink}（1） 开头（目的）}\par\vspace{1pt}
\begin{list}{}{\leftmargin=1.3em\labelwidth=0.95em\labelsep=0.35em\itemindent=0pt\topsep=0pt\partopsep=0pt\parsep=1.5pt\itemsep=0pt}
\item[{\color{accent}\scriptsize$\blacktriangleright$}]{\footnotesize In order to {\color{accent}\bfseries 【目的】}, {\color{accent}\bfseries 【主办方】} is to hold {\color{accent}\bfseries 【活动】} on {\color{accent}\bfseries 【日期】}.}\par\vspace{0.5pt}{\footnotesize\color{gray}in order to 置于句首是通知的标准开法，is to hold 表计划。}
\item[{\color{accent}\scriptsize$\blacktriangleright$}]{\footnotesize {\color{accent}\bfseries 【主办方】} will organize {\color{accent}\bfseries 【活动】} next {\color{accent}\bfseries 【时间】}, and all students are welcome.}\par\vspace{0.5pt}{\footnotesize\color{gray}主语用单位，不必写 I；all students are welcome 一句到位。}
\end{list}
\vspace{2pt}{\footnotesize\bfseries\color{ink}（2） 具体安排}\par\vspace{1pt}
\begin{list}{}{\leftmargin=1.3em\labelwidth=0.95em\labelsep=0.35em\itemindent=0pt\topsep=0pt\partopsep=0pt\parsep=1.5pt\itemsep=0pt}
\item[{\color{accent}\scriptsize$\blacktriangleright$}]{\footnotesize The activity will start at {\color{accent}\bfseries 【时间】} and last until {\color{accent}\bfseries 【时间】}.}\par\vspace{0.5pt}{\footnotesize\color{gray}起止时间都要写，只写开始时间属于要点不全。}
\item[{\color{accent}\scriptsize$\blacktriangleright$}]{\footnotesize Participants will {\color{accent}\bfseries 【内容一】} and then {\color{accent}\bfseries 【内容二】}.}\par\vspace{0.5pt}{\footnotesize\color{gray}用 participants 作主语，比 you will 更适合通知。}
\end{list}
\vspace{2pt}{\footnotesize\bfseries\color{ink}（3） 号召 / 联系}\par\vspace{1pt}
\begin{list}{}{\leftmargin=1.3em\labelwidth=0.95em\labelsep=0.35em\itemindent=0pt\topsep=0pt\partopsep=0pt\parsep=1.5pt\itemsep=0pt}
\item[{\color{accent}\scriptsize$\blacktriangleright$}]{\footnotesize Everyone is welcome, and your support will make a difference.}\par\vspace{0.5pt}{\footnotesize\color{gray}号召句要给出理由，光喊 Please come 没有力量。}
\item[{\color{accent}\scriptsize$\blacktriangleright$}]{\footnotesize For further information, please contact {\color{accent}\bfseries 【人】} at {\color{accent}\bfseries 【电话/邮箱】}.}\par\vspace{0.5pt}{\footnotesize\color{gray}通知必须留联系方式，否则格式与任务都不完整。}
\end{list}
\vspace{1pt}{\footnotesize\color{gray}语气与落款：单位与日期分两行，写在右下角；日期用 月 日, 年 的英语顺序。}\par
\endgroup
\vspace{4pt}\noindent{\color{accent}\bfseries 四、满分范文}\par\nopagebreak
{\color{rule}\rule{\linewidth}{0.6pt}}\par\nopagebreak
\vspace{1pt}
\vspace{2pt}\noindent\colorbox{soft}{\parbox{\dimexpr\linewidth-2\fboxsep\relax}{%
\vspace{2pt}
{\footnotesize\bfseries\color{accent}2013 真题话题：为贫困儿童举办慈善义卖}\hfill{\footnotesize\color{gray}正文 94 词}
\par\vspace{2pt}
\centerline{\footnotesize\bfseries Notice}\par\vspace{2pt}
\footnotesize In order to raise money for children in poor areas, the Student Union is to hold a charity sale on the playground this Saturday.\par
\footnotesize The sale will start at 9 a.m. and last until 4 p.m. Students and teachers are welcome to donate books, clothes and handmade gifts, or simply come and buy something. All the money raised will be sent to the children before the end of this month. Those who wish to donate should bring their items to Room 203 before Friday.\par
\footnotesize Everyone is welcome. For further information, please contact Li Ming at 12345678.\par
\par\noindent\hfill\footnotesize The Student Union\par
\par\noindent\hfill\footnotesize Oct. 20, 2025\par
\par\vspace{2pt}{\footnotesize\color{gray}通知：为贫困地区儿童筹款，学生会将于本周六在操场举办慈善义卖。义卖上午九点开始，持续到下午四点。欢迎师生捐赠书籍、衣物和手工艺品，也欢迎直接前来购买。所有筹得款项将于本月底前送到孩子们手中。有意捐赠者请在周五前将物品送到 203 室。欢迎大家参加。详情请联系李明，电话 12345678。学生会 2025 年 10 月 20 日}
\par\vspace{1pt}}}\par
\vspace{2pt}
\vspace{4pt}\noindent{\color{accent}\bfseries 五、雷区与得分点}\par\nopagebreak
{\color{rule}\rule{\linewidth}{0.6pt}}\par\nopagebreak
\vspace{1pt}
\vspace{1pt}\begingroup\setlength{\parindent}{0pt}
\begin{list}{}{\leftmargin=1.3em\labelwidth=0.95em\labelsep=0.35em\itemindent=0pt\topsep=0pt\partopsep=0pt\parsep=1.5pt\itemsep=0pt}
\item[{\color{warn}\scriptsize$\blacksquare$}]{\footnotesize 漏写标题 Notice 或落款，格式分直接损失，是通知类最常见错误。}
\item[{\color{warn}\scriptsize$\blacksquare$}]{\footnotesize 用 Dear students, 开头——通知不是书信，不需要称呼。}
\item[{\color{warn}\scriptsize$\blacksquare$}]{\footnotesize 只写“欢迎大家参加”，没有报名方式或联系方式，任务不完整。}
\item[{\color{warn}\scriptsize$\blacksquare$}]{\footnotesize 得分点：In order to 说明目的 + 起止时间与地点 + those who wish to 的报名要求 + 落款单位与日期。}
\end{list}
\endgroup
\clearpage
\label{g10}
\noindent\colorbox{accent}{%
\parbox{\dimexpr\linewidth-2\fboxsep\relax}{%
\color{white}\vspace{1pt}
{\LARGE\bfseries 10}\hspace{0.5em}{\Large 汇报 · 告知信}\hfill{\large\itshape Report / Letter of Information}
\\[1pt]{\footnotesize 真题 2019 · 2024}\par\vspace{1pt}}}
\par
\vspace{4pt}\noindent\colorbox{soft}{%
\parbox{\dimexpr\linewidth-2\fboxsep\relax}{\footnotesize{\color{accent}\bfseries 识别信号}\quad 题干出现 report / tell sb. about / 汇报计划、进展或情况，并常要求“征求意见”。}}
\par
\vspace{2pt}
\vspace{4pt}\noindent{\color{accent}\bfseries 一、任务要点}\par\nopagebreak
{\color{rule}\rule{\linewidth}{0.6pt}}\par\nopagebreak
\vspace{1pt}
\vspace{2pt}\noindent\begin{tabularx}{\dimexpr\linewidth-1.4pt\relax}{@{}>{\raggedright\arraybackslash}p{2.1cm}X@{}}
\toprule
\footnotesize{\color{accent}\bfseries 写作身份} & \footnotesize 项目负责人或学生代表，向老师、上级汇报 \\
\midrule
\footnotesize{\color{accent}\bfseries 必须写到} & \footnotesize （1）汇报什么（计划/进展）　（2）已完成与下一步　（3）具体安排（时间、分工）　（4）请对方在某日期前提意见 \\
\midrule
\footnotesize{\color{accent}\bfseries 一句话策略} & \footnotesize 汇报信靠“时间节点 + 分工”支撑可信度，一段里至少出现一个日期 \\
\bottomrule
\end{tabularx}\par
\vspace{2pt}
\vspace{4pt}\noindent{\color{accent}\bfseries 二、结构骨架}\par\nopagebreak
{\color{rule}\rule{\linewidth}{0.6pt}}\par\nopagebreak
\vspace{1pt}
\vspace{2pt}\noindent\begingroup\footnotesize\renewcommand{\arraystretch}{1.15}
\begin{tabularx}{\dimexpr\linewidth-1.4pt\relax}{@{}>{\raggedright\arraybackslash}p{2.4cm}>{\raggedright\arraybackslash}p{3.4cm}>{\raggedright\arraybackslash}X@{}}
\toprule
{\color{accent}\bfseries 位置} & {\color{accent}\bfseries 功能} & {\color{accent}\bfseries 必须落到纸面的内容} \\
\midrule
称呼 & 对师长 / 上级 & Dear Professor Wang, / Dear Mr. Smith, \\
第一段 & 点题 & 按上次商定，汇报……（一句话交代背景） \\
第二段 & 进展 & 已完成 + 下一步 + 时间节点或分工 \\
第三段 & 征求意见 & 请对方在某日期前就某方面给意见 \\
落款 & 格式 & Yours sincerely, + Li Ming \\
\bottomrule
\end{tabularx}\endgroup\par
\vspace{2pt}
\vspace{4pt}\noindent{\color{accent}\bfseries 三、可替换骨架句}\par\nopagebreak
{\color{rule}\rule{\linewidth}{0.6pt}}\par\nopagebreak
\vspace{1pt}
\vspace{1pt}\begingroup\setlength{\parindent}{0pt}
\vspace{2pt}{\footnotesize\bfseries\color{ink}（1） 点题}\par\vspace{1pt}
\begin{list}{}{\leftmargin=1.3em\labelwidth=0.95em\labelsep=0.35em\itemindent=0pt\topsep=0pt\partopsep=0pt\parsep=1.5pt\itemsep=0pt}
\item[{\color{accent}\scriptsize$\blacktriangleright$}]{\footnotesize I am writing to report on {\color{accent}\bfseries 【事项】}, as we agreed at our last meeting.}\par\vspace{0.5pt}{\footnotesize\color{gray}as we agreed 一句交代背景，省掉一整段铺垫。}
\item[{\color{accent}\scriptsize$\blacktriangleright$}]{\footnotesize Here is a brief account of {\color{accent}\bfseries 【事项】} so far.}\par\vspace{0.5pt}{\footnotesize\color{gray}account 比 report 轻，适合对熟人上级，口语信可用。}
\end{list}
\vspace{2pt}{\footnotesize\bfseries\color{ink}（2） 进展}\par\vspace{1pt}
\begin{list}{}{\leftmargin=1.3em\labelwidth=0.95em\labelsep=0.35em\itemindent=0pt\topsep=0pt\partopsep=0pt\parsep=1.5pt\itemsep=0pt}
\item[{\color{accent}\scriptsize$\blacktriangleright$}]{\footnotesize So far, we have {\color{accent}\bfseries 【已完成】}.}\par\vspace{0.5pt}{\footnotesize\color{gray}现在完成时交代进度，是汇报信的主干时态。}
\item[{\color{accent}\scriptsize$\blacktriangleright$}]{\footnotesize For the next stage, we plan to {\color{accent}\bfseries 【下一步】}.}\par\vspace{0.5pt}{\footnotesize\color{gray}plan to 后接动作，把“下一步”写实，而不是 we will continue。}
\end{list}
\vspace{2pt}{\footnotesize\bfseries\color{ink}（3） 征求意见}\par\vspace{1pt}
\begin{list}{}{\leftmargin=1.3em\labelwidth=0.95em\labelsep=0.35em\itemindent=0pt\topsep=0pt\partopsep=0pt\parsep=1.5pt\itemsep=0pt}
\item[{\color{accent}\scriptsize$\blacktriangleright$}]{\footnotesize I would appreciate your comments on {\color{accent}\bfseries 【方面】} before {\color{accent}\bfseries 【日期】}.}\par\vspace{0.5pt}{\footnotesize\color{gray}comments on 是征求书面意见的标准搭配。}
\item[{\color{accent}\scriptsize$\blacktriangleright$}]{\footnotesize Please let me know if any change is needed.}\par\vspace{0.5pt}{\footnotesize\color{gray}把决定权交回上级，语气得体。}
\end{list}
\vspace{1pt}{\footnotesize\color{gray}语气与落款：对老师、上级可用，一封信里用一次即可，多了显谄媚。}\par
\endgroup
\vspace{4pt}\noindent{\color{accent}\bfseries 四、满分范文}\par\nopagebreak
{\color{rule}\rule{\linewidth}{0.6pt}}\par\nopagebreak
\vspace{1pt}
\vspace{2pt}\noindent\colorbox{soft}{\parbox{\dimexpr\linewidth-2\fboxsep\relax}{%
\vspace{2pt}
{\footnotesize\bfseries\color{accent}2024 真题话题：古镇老宅保护调研计划}\hfill{\footnotesize\color{gray}正文 101 词}
\par\vspace{2pt}
\footnotesize Dear Professor Wang,\par
\footnotesize I am writing to report on our plan for the survey of the old houses in the ancient town, as we agreed at our last meeting.\par
\footnotesize We have decided to focus on twenty houses built before 1950 and to interview the families living in them. Two students will take photos and record the condition of the roofs and walls, while the rest will collect stories about the buildings. The field work will take three days at the end of this month, and we will submit a short report in early December.\par
\footnotesize I would appreciate your comments on our interview questions before 20 November.\par
\footnotesize Yours sincerely,\par
\footnotesize Li Ming\par
\par\vspace{2pt}{\footnotesize\color{gray}尊敬的王教授：按上次会议的商定，我向您汇报古镇老宅调研的计划。我们决定以二十座 1950 年前建成的老宅为对象，采访住在其中的家庭。两名同学负责拍照并记录屋顶和墙体的状况，其余同学收集关于这些建筑的故事。实地工作将在本月底进行三天，十二月初提交一份简短报告。如您能在 11 月 20 日前对我们的访谈问题提出意见，我将不胜感激。您真诚的，李明}
\par\vspace{1pt}}}\par
\vspace{2pt}
\vspace{4pt}\noindent{\color{accent}\bfseries 五、雷区与得分点}\par\nopagebreak
{\color{rule}\rule{\linewidth}{0.6pt}}\par\nopagebreak
\vspace{1pt}
\vspace{1pt}\begingroup\setlength{\parindent}{0pt}
\begin{list}{}{\leftmargin=1.3em\labelwidth=0.95em\labelsep=0.35em\itemindent=0pt\topsep=0pt\partopsep=0pt\parsep=1.5pt\itemsep=0pt}
\item[{\color{warn}\scriptsize$\blacksquare$}]{\footnotesize 只写“计划进展顺利”，没有时间节点和分工，内容空，落在 6 分以下。}
\item[{\color{warn}\scriptsize$\blacksquare$}]{\footnotesize 把汇报信写成建议信（反过来向对方提建议），身份颠倒，2019/2024 常见错。}
\item[{\color{warn}\scriptsize$\blacksquare$}]{\footnotesize 漏写“征求意见”这一步，任务要求未完成，最多 7 分。}
\item[{\color{warn}\scriptsize$\blacksquare$}]{\footnotesize 得分点：as we agreed 交代背景 + 现在完成时写已完成 + 具体分工与日期 + 明确征求哪方面的意见。}
\end{list}
\endgroup
\clearpage
\label{g11}
\noindent\colorbox{accent}{%
\parbox{\dimexpr\linewidth-2\fboxsep\relax}{%
\color{white}\vspace{1pt}
{\LARGE\bfseries 11}\hspace{0.5em}{\Large 推荐信}\hfill{\large\itshape Letter of Recommendation}
\\[1pt]{\footnotesize 大纲覆盖 · 模拟题高频（英语二真题未单独考）}\par\vspace{1pt}}}
\par
\vspace{4pt}\noindent\colorbox{soft}{%
\parbox{\dimexpr\linewidth-2\fboxsep\relax}{\footnotesize{\color{accent}\bfseries 识别信号}\quad 题干出现 recommend / 推荐某人做某事，或推荐一本书、一个地方；重点是“凭什么推荐”。}}
\par
\vspace{2pt}
\vspace{4pt}\noindent{\color{accent}\bfseries 一、任务要点}\par\nopagebreak
{\color{rule}\rule{\linewidth}{0.6pt}}\par\nopagebreak
\vspace{1pt}
\vspace{2pt}\noindent\begin{tabularx}{\dimexpr\linewidth-1.4pt\relax}{@{}>{\raggedright\arraybackslash}p{2.1cm}X@{}}
\toprule
\footnotesize{\color{accent}\bfseries 写作身份} & \footnotesize 老师、同学或朋友；先交代与被推荐人的关系与时长 \\
\midrule
\footnotesize{\color{accent}\bfseries 必须写到} & \footnotesize （1）推荐谁/什么、做什么　（2）两三条具体理由（经历 + 品质，用一件事证明）　（3）明确表态愿进一步说明 \\
\midrule
\footnotesize{\color{accent}\bfseries 一句话策略} & \footnotesize 推荐信不写形容词堆砌，只写“他做过什么”。一件事顶三个 excellent \\
\bottomrule
\end{tabularx}\par
\vspace{2pt}
\vspace{4pt}\noindent{\color{accent}\bfseries 二、结构骨架}\par\nopagebreak
{\color{rule}\rule{\linewidth}{0.6pt}}\par\nopagebreak
\vspace{1pt}
\vspace{2pt}\noindent\begingroup\footnotesize\renewcommand{\arraystretch}{1.15}
\begin{tabularx}{\dimexpr\linewidth-1.4pt\relax}{@{}>{\raggedright\arraybackslash}p{2.4cm}>{\raggedright\arraybackslash}p{3.4cm}>{\raggedright\arraybackslash}X@{}}
\toprule
{\color{accent}\bfseries 位置} & {\color{accent}\bfseries 功能} & {\color{accent}\bfseries 必须落到纸面的内容} \\
\midrule
称呼 & 不知姓名 / 知道姓名 & Dear Sir or Madam, / Dear Professor Wang, \\
第一段 & 点题 + 关系 & 我推荐谁 + 认识多久 + 推荐去做什么 \\
第二段 & 理由 & 职务或经历 + 一个具体成果 + 一项品质及其事例 \\
第三段 & 表态 & 毫不保留地推荐 + 可随时联系我 \\
落款 & 格式 & Yours sincerely, + Li Ming \\
\bottomrule
\end{tabularx}\endgroup\par
\vspace{2pt}
\vspace{4pt}\noindent{\color{accent}\bfseries 三、可替换骨架句}\par\nopagebreak
{\color{rule}\rule{\linewidth}{0.6pt}}\par\nopagebreak
\vspace{1pt}
\vspace{1pt}\begingroup\setlength{\parindent}{0pt}
\vspace{2pt}{\footnotesize\bfseries\color{ink}（1） 点题}\par\vspace{1pt}
\begin{list}{}{\leftmargin=1.3em\labelwidth=0.95em\labelsep=0.35em\itemindent=0pt\topsep=0pt\partopsep=0pt\parsep=1.5pt\itemsep=0pt}
\item[{\color{accent}\scriptsize$\blacktriangleright$}]{\footnotesize It is my pleasure to recommend {\color{accent}\bfseries 【人】} for {\color{accent}\bfseries 【职位/机会】}.}\par\vspace{0.5pt}{\footnotesize\color{gray}推荐信的标准开场，正式且不啰嗦。}
\item[{\color{accent}\scriptsize$\blacktriangleright$}]{\footnotesize I am writing to recommend {\color{accent}\bfseries 【书/地点】} to you, as you asked.}\par\vspace{0.5pt}{\footnotesize\color{gray}推荐物时用 to you；推荐人时用 for the position，介词不能混。}
\end{list}
\vspace{2pt}{\footnotesize\bfseries\color{ink}（2） 理由}\par\vspace{1pt}
\begin{list}{}{\leftmargin=1.3em\labelwidth=0.95em\labelsep=0.35em\itemindent=0pt\topsep=0pt\partopsep=0pt\parsep=1.5pt\itemsep=0pt}
\item[{\color{accent}\scriptsize$\blacktriangleright$}]{\footnotesize He has been {\color{accent}\bfseries 【职务】} for two years, during which he {\color{accent}\bfseries 【具体成绩】}.}\par\vspace{0.5pt}{\footnotesize\color{gray}职务 + 成绩，一句顶三句形容词。}
\item[{\color{accent}\scriptsize$\blacktriangleright$}]{\footnotesize He is patient with {\color{accent}\bfseries 【对象】} and quick at {\color{accent}\bfseries 【能力】}.}\par\vspace{0.5pt}{\footnotesize\color{gray}品质要落到“对谁有耐心、擅长什么”，否则等于没说。}
\end{list}
\vspace{2pt}{\footnotesize\bfseries\color{ink}（3） 表态}\par\vspace{1pt}
\begin{list}{}{\leftmargin=1.3em\labelwidth=0.95em\labelsep=0.35em\itemindent=0pt\topsep=0pt\partopsep=0pt\parsep=1.5pt\itemsep=0pt}
\item[{\color{accent}\scriptsize$\blacktriangleright$}]{\footnotesize I recommend him without reservation.}\par\vspace{0.5pt}{\footnotesize\color{gray}without reservation 意为“毫无保留”，推荐信固定表达。}
\item[{\color{accent}\scriptsize$\blacktriangleright$}]{\footnotesize Please feel free to contact me if you need further information.}\par\vspace{0.5pt}{\footnotesize\color{gray}留下可核实的姿态，增强可信度。}
\end{list}
\vspace{1pt}{\footnotesize\color{gray}语气与落款：表态句用 I am confident that，比 I think he is good 有分量。}\par
\endgroup
\vspace{4pt}\noindent{\color{accent}\bfseries 四、满分范文}\par\nopagebreak
{\color{rule}\rule{\linewidth}{0.6pt}}\par\nopagebreak
\vspace{1pt}
\vspace{2pt}\noindent\colorbox{soft}{\parbox{\dimexpr\linewidth-2\fboxsep\relax}{%
\vspace{2pt}
{\footnotesize\bfseries\color{accent}模拟题话题：推荐同学担任夏令营志愿者}\hfill{\footnotesize\color{gray}正文 97 词}
\par\vspace{2pt}
\footnotesize Dear Sir or Madam,\par
\footnotesize It is my pleasure to recommend my classmate Zhang Wei for the position of volunteer at your summer camp.\par
\footnotesize Zhang Wei has been chairman of our English Club for two years, during which he organized eight activities, including a speech contest for over 200 students. He is patient with children and quick at solving problems: last summer, when a visiting teacher fell ill, he rearranged the whole schedule within one afternoon.\par
\footnotesize I recommend him without reservation. I am sure he will be an asset to your team. Please feel free to contact me if you need further information.\par
\footnotesize Yours sincerely,\par
\footnotesize Li Ming\par
\par\vspace{2pt}{\footnotesize\color{gray}尊敬的先生/女士：我很荣幸推荐我的同学张伟担任贵夏令营的志愿者。张伟担任我校英语俱乐部主席两年，期间组织了八项活动，包括一场两百多名学生参加的演讲比赛。他对孩子有耐心，处理问题也很快：去年夏天一位来访教师突然生病，他一个下午就重新安排了全部日程。我毫无保留地推荐他。我相信他会成为贵团队的一份助力。如需进一步了解，请随时与我联系。您真诚的，李明}
\par\vspace{1pt}}}\par
\vspace{2pt}
\vspace{4pt}\noindent{\color{accent}\bfseries 五、雷区与得分点}\par\nopagebreak
{\color{rule}\rule{\linewidth}{0.6pt}}\par\nopagebreak
\vspace{1pt}
\vspace{1pt}\begingroup\setlength{\parindent}{0pt}
\begin{list}{}{\leftmargin=1.3em\labelwidth=0.95em\labelsep=0.35em\itemindent=0pt\topsep=0pt\partopsep=0pt\parsep=1.5pt\itemsep=0pt}
\item[{\color{warn}\scriptsize$\blacksquare$}]{\footnotesize 通篇 good, excellent, hard-working，没有一个事例，推荐没有说服力。}
\item[{\color{warn}\scriptsize$\blacksquare$}]{\footnotesize 不交代与被推荐人的关系，推荐资格受质疑，8 分以上必须写。}
\item[{\color{warn}\scriptsize$\blacksquare$}]{\footnotesize 推荐书或地点时照抄“推荐人”的模板（He is…），人称错，直接暴露背模板。}
\item[{\color{warn}\scriptsize$\blacksquare$}]{\footnotesize 得分点：关系时长 + 职务与成果 + 一件具体事例 + without reservation 的明确表态。}
\end{list}
\endgroup
\clearpage
\label{g12}
\noindent\colorbox{accent}{%
\parbox{\dimexpr\linewidth-2\fboxsep\relax}{%
\color{white}\vspace{1pt}
{\LARGE\bfseries 12}\hspace{0.5em}{\Large 申请 · 求职信}\hfill{\large\itshape Letter of Application}
\\[1pt]{\footnotesize 大纲覆盖 · 模拟题高频}\par\vspace{1pt}}}
\par
\vspace{4pt}\noindent\colorbox{soft}{%
\parbox{\dimexpr\linewidth-2\fboxsep\relax}{\footnotesize{\color{accent}\bfseries 识别信号}\quad 题干出现 apply for / 申请……岗位、志愿者或奖学金；常要求“说明资质”并“请求面试”。}}
\par
\vspace{2pt}
\vspace{4pt}\noindent{\color{accent}\bfseries 一、任务要点}\par\nopagebreak
{\color{rule}\rule{\linewidth}{0.6pt}}\par\nopagebreak
\vspace{1pt}
\vspace{2pt}\noindent\begin{tabularx}{\dimexpr\linewidth-1.4pt\relax}{@{}>{\raggedright\arraybackslash}p{2.1cm}X@{}}
\toprule
\footnotesize{\color{accent}\bfseries 写作身份} & \footnotesize 申请者；学历、经历写实，不夸大成“专家” \\
\midrule
\footnotesize{\color{accent}\bfseries 必须写到} & \footnotesize （1）申请什么、从何渠道得知　（2）匹配的资质与一个可量化成果　（3）请求面试 / 说明附件 \\
\midrule
\footnotesize{\color{accent}\bfseries 一句话策略} & \footnotesize 招聘方只关心两件事：你会做什么、什么时候能来。两句话都要出现 \\
\bottomrule
\end{tabularx}\par
\vspace{2pt}
\vspace{4pt}\noindent{\color{accent}\bfseries 二、结构骨架}\par\nopagebreak
{\color{rule}\rule{\linewidth}{0.6pt}}\par\nopagebreak
\vspace{1pt}
\vspace{2pt}\noindent\begingroup\footnotesize\renewcommand{\arraystretch}{1.15}
\begin{tabularx}{\dimexpr\linewidth-1.4pt\relax}{@{}>{\raggedright\arraybackslash}p{2.4cm}>{\raggedright\arraybackslash}p{3.4cm}>{\raggedright\arraybackslash}X@{}}
\toprule
{\color{accent}\bfseries 位置} & {\color{accent}\bfseries 功能} & {\color{accent}\bfseries 必须落到纸面的内容} \\
\midrule
称呼 & 不知姓名时 & Dear Sir or Madam, \\
第一段 & 点题 & 申请什么职位 + 从哪看到招聘信息 \\
第二段 & 资质 & 专业/学历 + 一段相关经历 + 一个可量化成果 + 可工作时间 \\
第三段 & 请求 & 请求面试 + 附简历 + 感谢 \\
落款 & 格式 & Yours sincerely, + Li Ming \\
\bottomrule
\end{tabularx}\endgroup\par
\vspace{2pt}
\vspace{4pt}\noindent{\color{accent}\bfseries 三、可替换骨架句}\par\nopagebreak
{\color{rule}\rule{\linewidth}{0.6pt}}\par\nopagebreak
\vspace{1pt}
\vspace{1pt}\begingroup\setlength{\parindent}{0pt}
\vspace{2pt}{\footnotesize\bfseries\color{ink}（1） 点题}\par\vspace{1pt}
\begin{list}{}{\leftmargin=1.3em\labelwidth=0.95em\labelsep=0.35em\itemindent=0pt\topsep=0pt\partopsep=0pt\parsep=1.5pt\itemsep=0pt}
\item[{\color{accent}\scriptsize$\blacktriangleright$}]{\footnotesize I am writing to apply for the position of {\color{accent}\bfseries 【职位】} advertised {\color{accent}\bfseries 【渠道】}.}\par\vspace{0.5pt}{\footnotesize\color{gray}advertised on your website 交代来源，一句话完成两项任务。}
\item[{\color{accent}\scriptsize$\blacktriangleright$}]{\footnotesize I wish to be considered for {\color{accent}\bfseries 【机会】}.}\par\vspace{0.5pt}{\footnotesize\color{gray}正式、简短，适合奖学金、项目申请。}
\end{list}
\vspace{2pt}{\footnotesize\bfseries\color{ink}（2） 资质}\par\vspace{1pt}
\begin{list}{}{\leftmargin=1.3em\labelwidth=0.95em\labelsep=0.35em\itemindent=0pt\topsep=0pt\partopsep=0pt\parsep=1.5pt\itemsep=0pt}
\item[{\color{accent}\scriptsize$\blacktriangleright$}]{\footnotesize I am a senior student majoring in {\color{accent}\bfseries 【专业】}, and I have {\color{accent}\bfseries 【经历】}.}\par\vspace{0.5pt}{\footnotesize\color{gray}专业 + 经历，一句交代“我是谁”。}
\item[{\color{accent}\scriptsize$\blacktriangleright$}]{\footnotesize Last term I helped {\color{accent}\bfseries 【对象】} improve {\color{accent}\bfseries 【方面】}, and {\color{accent}\bfseries 【成果】}.}\par\vspace{0.5pt}{\footnotesize\color{gray}用一个可量化的成果代替 qualified、capable 这类空词。}
\end{list}
\vspace{2pt}{\footnotesize\bfseries\color{ink}（3） 请求}\par\vspace{1pt}
\begin{list}{}{\leftmargin=1.3em\labelwidth=0.95em\labelsep=0.35em\itemindent=0pt\topsep=0pt\partopsep=0pt\parsep=1.5pt\itemsep=0pt}
\item[{\color{accent}\scriptsize$\blacktriangleright$}]{\footnotesize I would be grateful if you could offer me an interview.}\par\vspace{0.5pt}{\footnotesize\color{gray}请求面试的标准句，比 Please give me a chance 得体。}
\item[{\color{accent}\scriptsize$\blacktriangleright$}]{\footnotesize My resume is enclosed, and I can provide references if needed.}\par\vspace{0.5pt}{\footnotesize\color{gray}附件与证明人一并说明，显得有准备。}
\end{list}
\vspace{1pt}{\footnotesize\color{gray}语气与落款：把经历与岗位需求连起来，是申请信里最值钱的一句。}\par
\endgroup
\vspace{4pt}\noindent{\color{accent}\bfseries 四、满分范文}\par\nopagebreak
{\color{rule}\rule{\linewidth}{0.6pt}}\par\nopagebreak
\vspace{1pt}
\vspace{2pt}\noindent\colorbox{soft}{\parbox{\dimexpr\linewidth-2\fboxsep\relax}{%
\vspace{2pt}
{\footnotesize\bfseries\color{accent}模拟题话题：申请英语助教岗位}\hfill{\footnotesize\color{gray}正文 93 词}
\par\vspace{2pt}
\footnotesize Dear Sir or Madam,\par
\footnotesize I am writing to apply for the position of teaching assistant advertised on your website.\par
\footnotesize I am a senior student majoring in English, and I have worked as a part-time tutor for two years. Last term I helped a group of ten pupils improve their writing, and eight of them passed the final test with high marks. I am familiar with office software and available on weekends.\par
\footnotesize I would be grateful if you could offer me an interview. I can start work from the beginning of next month, and my resume is enclosed.\par
\footnotesize Yours sincerely,\par
\footnotesize Li Ming\par
\par\vspace{2pt}{\footnotesize\color{gray}尊敬的先生/女士：我写信申请贵网站招聘的助教岗位。我是英语专业大四学生，已做兼职家教两年。上学期我帮助十名小学生提高写作，其中八人以高分通过期末考试。我熟悉办公软件，周末可以工作。如能给我面试机会，我将不胜感激。我下月初即可开始工作，随信附上简历。您真诚的，李明}
\par\vspace{1pt}}}\par
\vspace{2pt}
\vspace{4pt}\noindent{\color{accent}\bfseries 五、雷区与得分点}\par\nopagebreak
{\color{rule}\rule{\linewidth}{0.6pt}}\par\nopagebreak
\vspace{1pt}
\vspace{1pt}\begingroup\setlength{\parindent}{0pt}
\begin{list}{}{\leftmargin=1.3em\labelwidth=0.95em\labelsep=0.35em\itemindent=0pt\topsep=0pt\partopsep=0pt\parsep=1.5pt\itemsep=0pt}
\item[{\color{warn}\scriptsize$\blacksquare$}]{\footnotesize 只写 I am qualified / I am interested，不写专业与经历，内容空。}
\item[{\color{warn}\scriptsize$\blacksquare$}]{\footnotesize 不说明“如何得知招聘信息”，要点缺失。}
\item[{\color{warn}\scriptsize$\blacksquare$}]{\footnotesize 写 I want to learn something here（把岗位当跳板），动机不当，语气失分。}
\item[{\color{warn}\scriptsize$\blacksquare$}]{\footnotesize 得分点：职位与来源 + 专业与两年经历 + 可量化成果（十人中八人高分）+ 可工作时间 + 请求面试。}
\end{list}
\endgroup
\clearpage
\label{app:a}
{\color{accent}\bfseries\LARGE 附录 A · 低分句 → 修改句 → 为什么}\par\vspace{4pt}
{\color{accent}\rule{\linewidth}{1.2pt}}\par
\vspace{4pt}\noindent{\footnotesize 下面 10 组是最常见的学生写法。左边是“能看懂但拿不到 8 分”的句子，右边是同一意思的达标写法。批改时按“必须改 / 可选优化”标注。}\par\vspace{4pt}
\vspace{2pt}\noindent\begingroup\footnotesize\renewcommand{\arraystretch}{1.15}
\begin{tabularx}{\dimexpr\linewidth-1.4pt\relax}{@{}>{\raggedright\arraybackslash}p{6.2cm}>{\raggedright\arraybackslash}p{7.0cm}>{\raggedright\arraybackslash}X@{}}
\toprule
{\color{accent}\bfseries 原句（常见写法）} & {\color{accent}\bfseries 修改句} & {\color{accent}\bfseries 说明} \\
\midrule
I am a student. I want to invite you to join our activity. & I am writing to invite you to our English speech contest, which will be held next Friday. & 两句并成一句主从句；点题句必须带活动名与时间，[必须改]。 \\
I very like this activity. & I am really fond of this activity. / I enjoy this activity very much. & like 是及物动词，不能用 very 直接修饰；程度副词位置也错，[必须改]。 \\
Thank you for your help, I finished my research. & Thanks to your help, I was able to finish my research on time. & 两个完整句不能用逗号连接（逗号粘连）；thanks to 更贴合致谢语境，[必须改]。 \\
Please tell me the information about the camp. & Could you tell me when the camp begins and how much it costs? & information 太笼统，阅卷人看不到“问了什么”；改为两个具体问题，[必须改]。 \\
I am sorry for I can't come. & I am sorry that I cannot come. / I apologize for being unable to come. & sorry for 后接名词或动名词，不接句子；cannot 在正式信中不缩写，[必须改]。 \\
The activity will hold on next Monday. & The activity will be held next Monday. / We will hold the activity next Monday. & hold 是及物动词，活动不能自己 hold；用被动或换主语，[必须改]。 \\
I hope you can give me a reply as soon as possible. & I would appreciate it if you could reply at your earliest convenience. & 原句不算错，但 as soon as possible 偏催促；公务信改客气句式，[可选优化]。 \\
You should read more books and practise more. & It might be a good idea to read more widely and to write one passage a day. & 建议信忌用 should 命令对方；给出可执行的具体做法，[必须改]。 \\
There are many students take part in it. & Many students take part in it. / There are many students who take part in it. & there be 句型后不能直接跟动词原形，缺关系代词，[必须改]。 \\
I am looking forward to meet you. & I am looking forward to meeting you. & look forward to 中 to 是介词，后接动名词，[必须改]。 \\
\bottomrule
\end{tabularx}\endgroup\par
\vspace{2pt}
\clearpage
\label{app:b}
{\color{accent}\bfseries\LARGE 附录 B · 功能句速查（按动作取用）}\par\vspace{4pt}
{\color{accent}\rule{\linewidth}{1.2pt}}\par\vspace{6pt}
\noindent{\bfseries\footnotesize\color{accent}点题（第一段）}\par\vspace{1pt}
\begin{list}{}{\leftmargin=1.3em\labelwidth=0.95em\labelsep=0.35em\itemindent=0pt\topsep=0pt\parsep=2pt\itemsep=2pt}
\item[{\color{accent}\scriptsize$\blacktriangleright$}]{\small I am writing to invite you to… / to apply for… / to complain about…}
\item[{\color{accent}\scriptsize$\blacktriangleright$}]{\small I am writing in reply to your letter of May 10, in which you asked about…}
\item[{\color{accent}\scriptsize$\blacktriangleright$}]{\small I was delighted to learn that… / I regret to say that…}
\item[{\color{accent}\scriptsize$\blacktriangleright$}]{\small On behalf of the Student Union, I would like to…}
\end{list}\vspace{4pt}
\noindent{\bfseries\footnotesize\color{accent}展开（第二段）}\par\vspace{1pt}
\begin{list}{}{\leftmargin=1.3em\labelwidth=0.95em\labelsep=0.35em\itemindent=0pt\topsep=0pt\parsep=2pt\itemsep=2pt}
\item[{\color{accent}\scriptsize$\blacktriangleright$}]{\small The activity will start at 9 a.m. and last until 4 p.m. in Room 203.}
\item[{\color{accent}\scriptsize$\blacktriangleright$}]{\small For one thing… For another… / First… Besides… More importantly…}
\item[{\color{accent}\scriptsize$\blacktriangleright$}]{\small What impressed me most was… / What I would do in your position is…}
\item[{\color{accent}\scriptsize$\blacktriangleright$}]{\small Those who wish to join should sign up before Friday.}
\end{list}\vspace{4pt}
\noindent{\bfseries\footnotesize\color{accent}请求与催促}\par\vspace{1pt}
\begin{list}{}{\leftmargin=1.3em\labelwidth=0.95em\labelsep=0.35em\itemindent=0pt\topsep=0pt\parsep=2pt\itemsep=2pt}
\item[{\color{accent}\scriptsize$\blacktriangleright$}]{\small I would be grateful if you could…}
\item[{\color{accent}\scriptsize$\blacktriangleright$}]{\small I would appreciate it if you could reply at your earliest convenience.}
\item[{\color{accent}\scriptsize$\blacktriangleright$}]{\small Please let me know by Thursday whether you are available.}
\item[{\color{accent}\scriptsize$\blacktriangleright$}]{\small Could you please send me some details about…?}
\end{list}\vspace{4pt}
\noindent{\bfseries\footnotesize\color{accent}感谢与道歉}\par\vspace{1pt}
\begin{list}{}{\leftmargin=1.3em\labelwidth=0.95em\labelsep=0.35em\itemindent=0pt\topsep=0pt\parsep=2pt\itemsep=2pt}
\item[{\color{accent}\scriptsize$\blacktriangleright$}]{\small I am writing to express my sincere gratitude for…}
\item[{\color{accent}\scriptsize$\blacktriangleright$}]{\small Without your help, I could hardly have…}
\item[{\color{accent}\scriptsize$\blacktriangleright$}]{\small Please accept my sincere apologies for…}
\item[{\color{accent}\scriptsize$\blacktriangleright$}]{\small I am sorry for any inconvenience I have caused.}
\end{list}\vspace{4pt}
\noindent{\bfseries\footnotesize\color{accent}收尾}\par\vspace{1pt}
\begin{list}{}{\leftmargin=1.3em\labelwidth=0.95em\labelsep=0.35em\itemindent=0pt\topsep=0pt\parsep=2pt\itemsep=2pt}
\item[{\color{accent}\scriptsize$\blacktriangleright$}]{\small I look forward to your early reply. / I am looking forward to meeting you.}
\item[{\color{accent}\scriptsize$\blacktriangleright$}]{\small I hope you will find these suggestions helpful.}
\item[{\color{accent}\scriptsize$\blacktriangleright$}]{\small Once again, thank you for your kindness.}
\item[{\color{accent}\scriptsize$\blacktriangleright$}]{\small Please feel free to contact me if you need further information.}
\end{list}\vspace{4pt}
\noindent{\bfseries\footnotesize\color{accent}衔接与过渡}\par\vspace{1pt}
\begin{list}{}{\leftmargin=1.3em\labelwidth=0.95em\labelsep=0.35em\itemindent=0pt\topsep=0pt\parsep=2pt\itemsep=2pt}
\item[{\color{accent}\scriptsize$\blacktriangleright$}]{\small To be specific, … / As for…, … / As you suggested, …}
\item[{\color{accent}\scriptsize$\blacktriangleright$}]{\small Unfortunately, … / To make up for it, …}
\item[{\color{accent}\scriptsize$\blacktriangleright$}]{\small In addition, … / What is more, … / On the other hand, …}
\item[{\color{accent}\scriptsize$\blacktriangleright$}]{\small Therefore, … / As a result, … / In this way, …}
\end{list}\vspace{4pt}
\clearpage
\label{app:c}
{\color{accent}\bfseries\LARGE 附录 C · 交卷前 60 秒自检清单}\par\vspace{4pt}
{\color{accent}\rule{\linewidth}{1.2pt}}\par\vspace{5pt}
\noindent{\footnotesize\color{gray}逐条打勾；任何一项没法打勾，就回去改那一句，不要改整篇。}\par\vspace{4pt}
\noindent{\bfseries\small\color{accent}格式}\par\vspace{1pt}
\begin{list}{}{\leftmargin=1.4em\labelwidth=1.0em\labelsep=0.4em\itemindent=0pt\topsep=0pt\parsep=2pt\itemsep=2pt}
\item[$\square$]{\footnotesize 称呼顶格且有逗号；结束语与署名各占一行、顶格。}
\item[$\square$]{\footnotesize 通知类有居中标题 Notice，右下角有单位与日期。}
\item[$\square$]{\footnotesize 署名是题目指定姓名，不是自己名字，也不是 Li Ming 加学号。}
\end{list}\vspace{4pt}
\noindent{\bfseries\small\color{accent}要点}\par\vspace{1pt}
\begin{list}{}{\leftmargin=1.4em\labelwidth=1.0em\labelsep=0.4em\itemindent=0pt\topsep=0pt\parsep=2pt\itemsep=2pt}
\item[$\square$]{\footnotesize 题干每个提示点都能在正文找到对应句（时间、地点、原因、要求、联系方式）。}
\item[$\square$]{\footnotesize 若是合并题（感谢＋邀请、祝贺＋建议、汇报＋征求意见），两件事各有一段或明确一句。}
\item[$\square$]{\footnotesize 咨询信有 2—3 个具体问题；投诉信有明确要求与期限。}
\end{list}\vspace{4pt}
\noindent{\bfseries\small\color{accent}语言}\par\vspace{1pt}
\begin{list}{}{\leftmargin=1.4em\labelwidth=1.0em\labelsep=0.4em\itemindent=0pt\topsep=0pt\parsep=2pt\itemsep=2pt}
\item[$\square$]{\footnotesize 主谓一致：The activity will be held…（不是 will hold）。}
\item[$\square$]{\footnotesize 时态成组：邀请用将来时，感谢用过去时，汇报用现在完成时。}
\item[$\square$]{\footnotesize 搭配清单：congratulations on / grateful for / apologize for / look forward to doing / be interested in。}
\item[$\square$]{\footnotesize 没有逗号粘连（两个完整句之间必须有连词或分号）。}
\item[$\square$]{\footnotesize 同一个句型不连续出现三次（I will… / I hope…）。}
\end{list}\vspace{4pt}
\noindent{\bfseries\small\color{accent}词数}\par\vspace{1pt}
\begin{list}{}{\leftmargin=1.4em\labelwidth=1.0em\labelsep=0.4em\itemindent=0pt\topsep=0pt\parsep=2pt\itemsep=2pt}
\item[$\square$]{\footnotesize 正文 90—120 词；不足 80 词要补一个细节句，超过 150 词删掉重复的客套。}
\item[$\square$]{\footnotesize 不把活动背景、节日来历写成流水账——这类内容不加分，还挤掉要点。}
\end{list}\vspace{4pt}
\vspace{6pt}\noindent\colorbox{soft}{\parbox{\dimexpr\linewidth-2\fboxsep\relax}{%
\footnotesize{\color{accent}\bfseries 最后一句}\quad 小作文的目标不是“写得漂亮”，而是“一件事说清楚、格式不出错、读者不用回头猜”。达到这一点，8 分起。}}
\end{document}